%%%%%%%%%%%%%%%%%%%%%%%%%%%%%%%%%%%%%%%%%%%%%%%%%%%%%%%%%%%%%%%%%%%%%%%%%%%%%%
%  Marketing Management (IM345AT) -- Volume 2 of 5
%  UNIT II : Social Media Marketing and Search Engine Optimization
%%%%%%%%%%%%%%%%%%%%%%%%%%%%%%%%%%%%%%%%%%%%%%%%%%%%%%%%%%%%%%%%%%%%%%%%%%%%%%

\documentclass{MMTextbook}

\unitnumber{II}
\unittitle{Social Media Marketing and Search Engine Optimization}
\unitsubtitle{Platforms, audience insight, campaign optimisation, and the
complete search optimisation stack}
\unithours{08 Hours}
\volumeno{2}

\makeindex[intoc=true,columns=2,title=Index]

\hypersetup{
  pdftitle={Marketing Management IM345AT -- Unit II: Social Media Marketing and
             Search Engine Optimization},
  pdfsubject={Semester IV, Department of Industrial Engineering and Management,
              RV College of Engineering},
  pdfkeywords={social media marketing, SEO, SERP, backlinks, IM345AT,
               marketing management}
}

\begin{document}

%%%%%%%%%%%%%%%%%%%%%%%%%%%%%%%%%%%%%%%%%%%%%%%%%%%%%%%%%%%%%%%%%%%%%%%%%%%%%%
\frontmatter
\pagestyle{mmfront}

\halftitlepage
\maintitlepage

\colophonpage{%
\textbf{\coursetitle\ ---\ The Study Text}\\[2pt]
Volume 2 of 5 \textbullet\ Unit II: Social Media Marketing and Search Engine
Optimization\\[10pt]
%
Prepared for the course \coursetitle\ (\coursecode), Semester IV, of the
\coursedept, \courseinstitute.\\[10pt]
%
\textbf{Sources.} The syllabus, course outcomes, assessment rubrics and
reference-book list reproduced in the front matter are transcribed from the
official course document. The previous-year questions in Appendix~B are
transcribed from RVCE Continuous Internal Evaluation and Semester End
Examination papers for this course, and where an official scheme of valuation
was released the model answers follow it. The explanatory text is compiled from
classroom lecture notes and, where those were incomplete, from the prescribed
reference texts.\\[10pt]
%
\textbf{Typography.} Set in Nimbus Roman with URW Gothic display headings.
Composed in \LaTeX{} using the \texttt{MMTextbook} class written for this
series. The visual design is derived from \emph{The Legrand Orange Book}
template by Mathias Legrand, with modifications by Vel, distributed through
LaTeXTemplates.com under a Creative Commons BY-NC-SA 4.0 licence; the class
re-implements that design language on standard CTAN packages. All figures are
drawn in \TikZ{} and \texttt{pgfplots}.\\[10pt]
%
\textbf{Status.} Unofficial student-compiled study material. Not issued,
endorsed or verified by the Department or by RV College of Engineering.\\[10pt]
%
2026 edition.}

% ---------------------------------------------------------------- preface
\frontchapter{Preface to Volume 2}

Unit~II is the largest unit in the course --- eight hours against seven for each
of the others --- and it is the unit that carries the most machinery. It has two
halves that look separate and are not. Social media marketing distributes
content to people; search engine optimisation makes content findable by people.
Both are answers to the same question: how does a stranger who has never heard
of you come to trust you enough to buy?

The first half opens with a single statistic that governs everything after it.
Roughly nine people in ten believe what their peers say about a brand, while
fewer than half believe what the brand says about itself. Every construct in
Chapters~1 to~5 --- the lifecycle orbit, the 4-1-1 content ratio, the algorithm
filter, the campaign optimisation levers --- follows from that asymmetry. The
objective is never to talk louder. It is to build something the audience will
carry on your behalf.

The second half is more mechanical and, for that reason, more reliably scored.
Search optimisation divides into three layers that a well-known analogy calls a
house: technical work is the foundation that makes a site discoverable, on-page
work is the living floor that makes it relevant, and off-page work is the roof
that makes it authoritative. Rankings are a function of all three, and a
question that names only one has answered a third of it.

Chapters~1 to~5 cover social media marketing; Chapters~6 to~11 cover search
engine optimisation. The chapter sequence follows the syllabus phrase order
exactly, and Appendix~A records the mapping.

Two editorial notes specific to this volume.

First, several topics in this unit exist in \emph{two} taught versions --- the
customer lifecycle orbit and the 4~Rs of content scaling both appear one way in
the lecture material and another way in a released scheme of valuation. Where
that happens, this volume gives both, says which is which, and recommends the
answer that satisfies either marker. Those passages are flagged with a
\emph{Two Accepted Versions} panel.

Second, one topic in the source material --- Search Everywhere Optimization,
sometimes called SEO 2.0 --- is a forward-looking extension that no unit
descriptor names. It is set out in Appendix~C rather than in a numbered chapter,
for the same reason the supplementary topics of Volume~1 were: the numbered
chapters carry the syllabus and nothing else, but genuinely useful material
should not simply vanish.

\vspace{12pt}
\noindent\emph{The whole of the course syllabus is reproduced overleaf. The
unit covered by this volume is Unit~II.}

% ------------------------------------------------- conventions & syllabus
%%%%%%%%%%%%%%%%%%%%%%%%%%%%%%%%%%%%%%%%%%%%%%%%%%%%%%%%%%%%%%%%%%%%%%%%%%%%%%
%  mm-howtouse.tex
%  Editorial conventions for the series.  Shared, identical, across all five
%  volumes so that a reader who learns the apparatus once can use any volume.
%%%%%%%%%%%%%%%%%%%%%%%%%%%%%%%%%%%%%%%%%%%%%%%%%%%%%%%%%%%%%%%%%%%%%%%%%%%%%%

\frontchapter{How to Use This Study Text}

\section*{The shape of a volume}
\addcontentsline{toc}{section}{The shape of a volume}

Each of the five volumes in this series covers exactly one unit of the
IM345AT syllabus and is self-contained. Every volume opens with the official
syllabus for the whole course, so that the unit in hand can always be located
within it, and closes with a set of appendices carrying the concordance,
the question bank, the glossary and a quick-reference sheet.

The numbered chapters of each volume map one-to-one onto the topic phrases of
that unit's official syllabus descriptor. The order of the chapters is the
order of the syllabus. Appendix~A sets out that mapping explicitly, phrase by
phrase, so coverage can be audited rather than assumed.

\section*{The five panel types}
\addcontentsline{toc}{section}{The five panel types}

Material that is not continuous prose is set in one of five panels. Each has a
fixed meaning throughout the series.

\begin{definitionbox}[the term being defined]
A formal definition, worth reproducing close to verbatim in an answer script.
The term appears in the panel heading and again in the index.
\end{definitionbox}

\begin{keybox}[High Yield]
A result, list or distinction that carries disproportionate examination weight.
Where a topic has recurred across papers, the recurrence is recorded as a
bracketed tag. \pyq{CIE-I, 3 Apr 2025; CIE-I, 16 Apr 2026}
\end{keybox}

\begin{exambox}[Model Answer]
An answer written at the length and in the register the marking scheme
rewards. Where an official scheme of valuation was released, the panel follows
it.
\end{exambox}

\begin{examplebox}[Worked Example]
A concrete illustration --- a named company, a named persona, a worked
calculation --- rather than a general statement.
\end{examplebox}

\begin{warnbox}
A common and costly error: a distinction usually collapsed, a term usually
misused, or an answer that looks complete and is not.
\end{warnbox}

A sixth panel, in green, sets out a framework or process as a single
sequence:

\begin{frameworkbox}[Framework]
Step one \ra\ step two \ra\ step three \ra\ back to step one.
\end{frameworkbox}

\section*{How previous-year questions are recorded}
\addcontentsline{toc}{section}{How previous-year questions are recorded}

\begin{keybox}[The Bracket Convention]
A past-paper appearance is recorded \textbf{only} as a bracketed source tag
attached to the material, never as a sentence of narrative. So the text reads
\begin{quote}\sffamily\footnotesize
Push versus pull marketing. \pyq{CIE-I, 3 Apr 2025; CIE-I, 16 Apr 2026; SEE, Jun/Jul 2025 --- 2 M}
\end{quote}
and never ``this was asked in 2025 and again in 2026''. The tag carries the
paper, the date and, where the mark allocation is known, the marks. Multiple
appearances are separated by semicolons.
\end{keybox}

\noindent The abbreviations used inside tags are:

\vspace{4pt}
{\sffamily\small
\begin{tabular}{@{}P{0.14\textwidth}P{0.80\textwidth}@{}}
\toprule[1.1pt]
\rlab{CIE-I} & First Continuous Internal Evaluation test\\
\rlab{CIE-II} & Second Continuous Internal Evaluation test\\
\rlab{CIE-III} & Third Continuous Internal Evaluation test\\
\rlab{SEE} & Semester End Examination\\
\rlab{2 M} & Mark allocation of the question as printed on the paper\\
\bottomrule[1.1pt]
\end{tabular}}

\vspace{10pt}

Appendix~B of each volume reproduces every past-paper question traced to that
unit, in full, grouped by mark weight, each carrying the same bracketed tag.

\section*{Cross-references, figures and the index}
\addcontentsline{toc}{section}{Cross-references, figures and the index}

Cross-references are given as \S\,chapter.section --- so \S\,1.4 is the fourth
section of Chapter~1. Every figure in this series is drawn in \TeX{} rather
than imported as an image, so the diagrams print at full resolution at any
size and use the same typeface as the surrounding text. Figures and tables are
listed after the table of contents.

Terms set in \textbf{\textit{bold italic}} on first use are indexed. The index
at the back of each volume is a real back-of-book index built from those
terms, not a repeat of the contents list.

\section*{Status of this text}
\addcontentsline{toc}{section}{Status of this text}

\begin{warnbox}[Unofficial Document]
This is a student-compiled study text. It is not issued, endorsed or verified
by the Department of Industrial Engineering and Management or by RV College of
Engineering. The syllabus reproduced in the front matter, the assessment
rubrics and the previous-year questions are transcribed from official
documents; the explanatory prose around them is not official. Where the source
lecture notes were incomplete, gaps have been filled from the prescribed
reference texts listed in the front matter. Always cross-check against the
official syllabus and your own class notes.
\end{warnbox}

%%%%%%%%%%%%%%%%%%%%%%%%%%%%%%%%%%%%%%%%%%%%%%%%%%%%%%%%%%%%%%%%%%%%%%%%%%%%%%
%  mm-syllabus.tex
%  Verbatim reproduction of the official IM345AT course document.
%  Shared, byte-identical, across all five volumes of the series.
%%%%%%%%%%%%%%%%%%%%%%%%%%%%%%%%%%%%%%%%%%%%%%%%%%%%%%%%%%%%%%%%%%%%%%%%%%%%%%

\frontchapter{Official Syllabus: IM345AT}

\noindent
{\sffamily\small\color{slatelight}%
Reproduced from the official course document issued by the Department of
Industrial Engineering and Management, RV College of Engineering. The text of
the five unit descriptors below is the authority against which every chapter of
this series is written; nothing outside it appears in the numbered chapters of
any volume.\par}

\vspace{10pt}

\begin{center}
\sffamily\small
\begin{tabular}{@{}P{0.238\textwidth}P{0.238\textwidth}P{0.218\textwidth}P{0.218\textwidth}@{}}
\toprule[1.1pt]
\multicolumn{4}{@{}l}{\bfseries\normalsize Semester: IV \quad\textbullet\quad MARKETING MANAGEMENT}\\
\multicolumn{4}{@{}l}{\color{slatelight}Category: Professional Core Course (Theory)}\\
\midrule
\rlab{Course Code} & IM345AT & \rlab{CIE Marks} & 100 Marks\\
\rlab{Credits (L:T:P)} & 3:0:0 & \rlab{SEE Marks} & 100 Marks\\
\rlab{Total Hours} & 45 L & \rlab{SEE Duration} & 3.00 Hours\\
\bottomrule[1.1pt]
\end{tabular}
\end{center}

\vspace{6pt}

% ---------------------------------------------------------------- unit blocks
\newcommand{\syllunit}[3]{%
  \begin{tcolorbox}[mmbase, colback=white, colframe=ocre!45, boxrule=0.7pt,
      sharp corners, fontupper=\small,
      fonttitle=\sffamily\bfseries\small\color{white},
      coltitle=white, colbacktitle=#3,
      title={#1},
      attach boxed title to top left={xshift=-0.7pt, yshift=-0.7pt},
      boxed title style={sharp corners, boxrule=0pt, colback=#3},
      top=16pt]
    #2
  \end{tcolorbox}}

\syllunit{UNIT--I \hfill 07 Hrs}{%
\textbf{Introduction to Digital Marketing:} Principles of Digital Marketing;
Digital Marketing Channels; Tools to Create Buyer Persona; Competitor Research
Tools, Website Analysis Tools, etc.\par
\vspace{3pt}
\textbf{Content Marketing:} Content Marketing Concepts \& Strategies; Planning,
Creating, Distributing \& Promoting Content; Optimize Website UX \& Landing
Pages; Measure Impact; Metrics \& Performance; Using Content Research for
Opportunities, etc.}{ocre}

\syllunit{UNIT--II \hfill 08 Hrs}{%
\textbf{Social Media Marketing:} Introduction; Major Social Media Platforms for
Marketing; Developing Data-driven Audience \& Campaign Insights; Social Media
for Business; Creation \& Optimization of Social Media Campaigns, etc.\par
\vspace{3pt}
\textbf{Search Engine Optimization:} Search Engine Optimization Fundamentals;
Keywords and SEO Content Plan; SEO \& Business Objectives; Writing SEO Content;
On-site \& off-site SEO; Optimize Organic Search Ranking, etc.}{ocre}

\syllunit{UNIT--III \hfill 07 Hrs}{%
\textbf{Web Analytics \& Google Analytics:} Google Analytics Tools; Web
Analytics Tools, etc.\par
\vspace{3pt}
\textbf{E-mail Marketing:} Effective E-mail Campaigns; E-mail Plan; E-mail
Marketing Campaign Analysis; Measuring Conversions \& keeping up, etc.}{ocre}

\syllunit{UNIT--IV \hfill 07 Hrs}{%
\textbf{Web Design:} Web design, optimization of websites; Publishing a basic
website; User-centered Design and Website Optimization; Design Principles and
Website Copy; Website Metrics \& Developing Insight, etc.\par
\vspace{3pt}
\textbf{Mobile Marketing:} Difference between mobile advertising and marketing,
utilizing mobile marketing for sales promotions, online applications,
etc.}{ocre}

\syllunit{UNIT--V \hfill 07 Hrs}{%
\textbf{Conversion Optimization:} What is AIDAS and its role; website
optimization; what visitors want to see on the website; how to optimize key
element and increase the effect of landing on a particular page\par
\vspace{3pt}
\textbf{Digital Analytics:} Evolution of Digital Analytics, information about
end-to-end customer experience, analyst's influence on business, role as a
change agent, etc.}{ocre}

% ------------------------------------------------------------ course outcomes
\section*{Course Outcomes}
\addcontentsline{toc}{section}{Course Outcomes}

\noindent{\small After completing the course, the students will be able to:\par}
\vspace{4pt}

{\sffamily\small
\begin{longtable}{@{}P{0.07\textwidth}P{0.89\textwidth}@{}}
\toprule[1.1pt]
\rlab{CO1} & Differentiate the benefits drawn by updated marketing mix from
traditional marketing mix for effective marketing management there by to stay
competitive in today's global market-place.\\
\midrule
\rlab{CO2} & Develop an effective holistic marketing atmosphere to efficiently
face the challenges in dynamically changing market.\\
\midrule
\rlab{CO3} & Formulate a potential marketing plan to effectively reach the
targeted market segments, by delivering the value to targeted customers through
practicing sound marketing research.\\
\midrule
\rlab{CO4} & Create new channels to improvise marketing to achieve and maintain
competitive position in globalized market-place.\\
\bottomrule[1.1pt]
\end{longtable}}

% --------------------------------------------------------------- reference books
\section*{Prescribed Reference Books}
\addcontentsline{toc}{section}{Prescribed Reference Books}

{\small
\begin{enumerate}[leftmargin=2.1em]
\item \textit{Marketing Management}, Philip Kotler, Kevin Lane Keller, 15th
      Edition, 2016, Pearson, ISBN: 978-93-325-5718-5.
\item \textit{Digital Marketing --- Strategy, Implementation \& Practice}, Dave
      Chaffey, Fiona Ellis-Chadwick, 7th Edition, 2019, Pearson,
      ISBN: 9781292241623, 1292241624.
\item \textit{Marketing Research}, Donald S. Tull, Del I. Hawkins, 6th Edition,
      Prentice Hall India, ISBN: 8120309618.
\item \textit{Marketing Management --- A South Asian Perspective}, Philip
      Kotler, Kevin Lane Keller, Abraham Koshy, Mithileshwar Jha, 14th Edition,
      2013, Pearson, ISBN: 978-81-317-6716-0.
\item \textit{Marketing Research}, David A. Aaker, V. Kumar, George S. Day, 9th
      Edition, 2008, John Wiley \& Sons, ISBN: 978-265-1791-6.
\end{enumerate}}

% ------------------------------------------------------------------- rubrics
\frontchapter{Assessment Scheme}

\section*{Rubric for the Continuous Internal Evaluation (Theory)}
\addcontentsline{toc}{section}{Rubric for the Continuous Internal Evaluation}

{\sffamily\small
\begin{longtable}{@{}P{0.045\textwidth}P{0.80\textwidth}>{\raggedleft\arraybackslash}P{0.07\textwidth}@{}}
\toprule[1.1pt]
\rowcolor{ocre}\thd{\#} & \thd{Component} & \thd{Marks}\\
\midrule
\rlab{1.} & \textbf{Quizzes:} Quizzes will be conducted in online/offline mode.
TWO QUIZZES will be conducted \& each quiz will be evaluated for 10 marks. The
sum of two quizzes will be the final quiz marks. & \bfseries 20\\
\midrule
\rowcolor{tablealt}
\rlab{2.} & \textbf{Tests:} Students will be evaluated in test, descriptive
questions with different complexity levels (Revised Bloom's Taxonomy Levels:
Remembering, Understanding, Applying, Analyzing, Evaluating and Creating). TWO
tests will be conducted. Each test will be evaluated for 50 marks, adding up to
100 marks. Final test marks will be reduced to 40 marks. & \bfseries 40\\
\midrule
\rlab{3.} & \textbf{Experiential Learning:} Students will be evaluated for their
creativity and practical implementation of the problem. Case study-based
teaching learning (10), Program specific requirements (10), Video based
seminar/presentation/demonstration (20), adding up to 40. & \bfseries 40\\
\midrule
\multicolumn{2}{@{}l}{\bfseries MAXIMUM MARKS FOR THE CIE THEORY} & \bfseries 100\\
\bottomrule[1.1pt]
\end{longtable}}

\section*{Rubric for the Semester End Examination (Theory)}
\addcontentsline{toc}{section}{Rubric for the Semester End Examination}

{\sffamily\small
\begin{longtable}{@{}P{0.12\textwidth}P{0.72\textwidth}>{\raggedleft\arraybackslash}P{0.08\textwidth}@{}}
\toprule[1.1pt]
\rowcolor{ocre}\thd{Q. No.} & \thd{Contents} & \thd{Marks}\\
\midrule
\multicolumn{3}{@{}l}{\bfseries\color{slate} PART A}\\
\midrule
\rlab{1} & Objective type questions covering entire syllabus & \bfseries 20\\
\midrule
\multicolumn{3}{@{}l}{\bfseries\color{slate} PART B \quad\normalfont\itshape(Maximum of TWO sub-divisions only)}\\
\midrule
\rlab{2}      & Unit 1: (Compulsory)      & \bfseries 16\\
\rowcolor{tablealt}
\rlab{3 \& 4} & Unit 2: (Internal Choice) & \bfseries 16\\
\rlab{5 \& 6} & Unit 3: (Internal Choice) & \bfseries 16\\
\rowcolor{tablealt}
\rlab{7 \& 8} & Unit 4: (Internal Choice) & \bfseries 16\\
\rlab{9 \& 10}& Unit 5: (Internal Choice) & \bfseries 16\\
\midrule
\multicolumn{2}{@{}l}{\bfseries TOTAL} & \bfseries 100\\
\bottomrule[1.1pt]
\end{longtable}}

\vspace{8pt}

\begin{keybox}[What the SEE Structure Means for This Volume]
Part A draws twenty marks of objective questions from the \emph{entire}
syllabus, so no unit can be skipped. In Part B, Question~2 is compulsory and is
always set from Unit~I; Units~II to~V each carry a 16-mark internal-choice
pair. A candidate therefore answers Unit~I plus one question from each of the
remaining four units --- five full questions in all.
\end{keybox}


% ---------------------------------------------------------------- contents
\cleardoublepage
\pdfbookmark[0]{Contents}{toc}
\tableofcontents
\cleardoublepage
\listoffigures
\cleardoublepage
\listoftables

%%%%%%%%%%%%%%%%%%%%%%%%%%%%%%%%%%%%%%%%%%%%%%%%%%%%%%%%%%%%%%%%%%%%%%%%%%%%%%
\mainmatter
\pagestyle{mmmain}

%%%%%%%%%%%%%%%%%%%%%%%%%%%%%%%%%%%%%%%%%%%%%%%%%%%%%%%%%%%%%%%%%%%%%%%%%%%%%%
\chapter{Social Media Marketing --- Introduction}
%%%%%%%%%%%%%%%%%%%%%%%%%%%%%%%%%%%%%%%%%%%%%%%%%%%%%%%%%%%%%%%%%%%%%%%%%%%%%%

\syllabusstrip{Social Media Marketing: Introduction.}

\begin{learningoutcomes}
\begin{itemize}
\item What social media marketing is, and its primary goal.
\item The trust asymmetry, and why it dictates the whole strategy.
\item The omnichannel gearbox --- three ways social media drives the business.
\item Seven myths about social media marketing, and what is actually true.
\item The characteristics of a competent social media marketer.
\item The customer lifecycle orbit, and the economics of retention.
\end{itemize}
\end{learningoutcomes}

\section{Definition and Primary Goal}

\begin{definitionbox}[Social Media Marketing]
\keytermas{Social media marketing}{social media marketing} (SMM) is a form of
internet marketing that involves creating and sharing content on social media
networks in order to achieve marketing and branding goals. It includes
publishing text, images and video, running paid social advertising, and engaging
directly with the audience.
\end{definitionbox}

\begin{exambox}[Primary Goal --- 2 Marks]
\pyq{SEE, Jun/Jul 2025 --- 2 M}
The primary goal of social media marketing is to build brand awareness\ix{brand awareness},
engagement and community with a target audience on social platforms, and to
convert that engagement into traffic, leads, sales and long-term customer
loyalty.
\end{exambox}

\begin{warnbox}[Social Media Is Not a Series of Posts]
The most common conceptual error in this unit is treating social media marketing
as publishing. It is an engineered system with an audience model, a content
ratio, a distribution algorithm to satisfy, a measurement layer and an
optimisation loop. Every chapter that follows is a component of that system.
\end{warnbox}

\section{The Word-of-Mouth Imperative}

\begin{keybox}[The Trust Asymmetry]
Roughly \textbf{90\%} of people believe what their peers say about a brand,
while only about \textbf{48\%} trust what businesses say directly. Audiences
increasingly read brand-to-buyer communication as advertising, which means trust
is the indispensable connective tissue between a brand and customer loyalty.

\smallskip
People also value word-of-mouth recommendations from friends and family about
\textbf{41\%} more than social media recommendations --- so even inside social
media, the personal channel outperforms the branded one.
\end{keybox}

\begin{figure}[htbp]
\centering
\begin{tikzpicture}
\begin{axis}[
  width=11.6cm, height=4.3cm,
  font=\sffamily\footnotesize,
  xbar, y=0.82cm,
  xmin=0, xmax=112,
  xtick={0,20,40,60,80,100},
  xticklabel={\pgfmathprintnumber{\tick}\%},
  tick label style={font=\sffamily\scriptsize,text=slatelight},
  axis lines=left, axis line style={rulegrey},
  ytick=data,
  symbolic y coords={brand,peers},
  yticklabel style={font=\sffamily\small,text=slate,align=right},
  nodes near coords, nodes near coords align={horizontal},
  every node near coord/.append style={font=\sffamily\bfseries\scriptsize,
        text=slate},
  enlarge y limits=0.42]
  \addplot[draw=greenok,fill=greenok!55,bar width=0.42cm] coordinates {(90,peers)};
  \addplot[draw=warnred,fill=warnred!45,bar width=0.42cm] coordinates {(48,brand)};
\end{axis}
\node[anchor=west,font=\sffamily\scriptsize,text=slate,align=left]
  at (7.3,1.35) {believe what\\ \textbf{peers} say};
\node[anchor=west,font=\sffamily\scriptsize,text=slate,align=left]
  at (7.3,0.05) {trust what\\ \textbf{businesses} say};
\end{tikzpicture}
\caption[The trust asymmetry]%
        {The trust asymmetry. Because the peer channel is roughly twice as
         credible as the brand channel, the strategic objective is not to speak
         more loudly but to make the audience willing to speak for you.}
\label{fig:trustasymmetry}
\end{figure}

\begin{keybox}[The Strategic Conclusion]
The goal is not to talk \emph{at} the audience. It is to engineer an experience
valuable enough that they do the marketing \emph{for} you.
\end{keybox}

\section{The Omnichannel Gearbox}

Social media drives the business through three distinct mechanisms. Treating
them as one --- ``engagement'' --- is what makes social media budgets hard to
defend.

\begin{figure}[htbp]
\centering
\begin{tikzpicture}[font=\sffamily\footnotesize,
    gear/.style={draw=ocre,line width=1pt,fill=ocrepale,circle,
                 minimum size=2.65cm,align=center,
                 font=\sffamily\bfseries\scriptsize,text=slate}]
  \node[gear] (g1) at (0,0)    {SOCIAL\\ SELLING};
  \node[gear] (g2) at (3.5,0)  {SEO\\ FUEL};
  \node[gear] (g3) at (7.0,0)  {OMNICHANNEL\\ SYNC};

  \draw[-{Latex[length=2.3mm]},ocre,line width=1pt] (1.4,0.42) -- (2.1,0.42);
  \draw[-{Latex[length=2.3mm]},ocre,line width=1pt] (2.1,-0.42) -- (1.4,-0.42);
  \draw[-{Latex[length=2.3mm]},ocre,line width=1pt] (4.9,0.42) -- (5.6,0.42);
  \draw[-{Latex[length=2.3mm]},ocre,line width=1pt] (5.6,-0.42) -- (4.9,-0.42);

  \node[anchor=north,text=slate,font=\sffamily\scriptsize,align=center,
        text width=3.1cm] at (0,-1.55)
    {one-to-one sales connections; personal, relevant content straight to the
     prospect};
  \node[anchor=north,text=slate,font=\sffamily\scriptsize,align=center,
        text width=3.1cm] at (3.5,-1.55)
    {shares produce links and time-on-site, which lift organic ranking
     indirectly};
  \node[anchor=north,text=slate,font=\sffamily\scriptsize,align=center,
        text width=3.1cm] at (7.0,-1.55)
    {coordinated messaging becomes one continuous conversation, not silos};
\end{tikzpicture}
\caption[The omnichannel gearbox]%
        {The three gears through which social media marketing drives business
         outcomes. Each is measured differently, which is why they must be
         argued for separately.}
\label{fig:gearbox}
\end{figure}

\begin{enumerate}
\item \sfb{Social selling} enables one-to-one sales connections. Sales teams
      offer personal, relevant content directly to prospects, accelerating the
      purchase decision.
\item \sfb{SEO fuel.} Social shares push content outward. The resulting incoming
      links and increased time-on-site engagement communicate high value to
      search engines, indirectly lifting organic rankings.
\item \sfb{Omnichannel synchronization.} Coordinated messaging across channels
      creates one continuous conversation, preventing siloed and contradictory
      customer experiences.
\end{enumerate}

\section{Seven Myths}

\begin{mmlongtable}{Myths about social media marketing}{tab:myths}%
{@{}P{0.31\textwidth}P{0.62\textwidth}@{}}
\rowcolor{ocre}\thd{The myth} & \thd{What is actually true}\\
\midrule
\endfirsthead
\rowcolor{ocre}\thd{The myth} & \thd{What is actually true}\\
\midrule
\endhead
\rlab{``Social media is free.''} & The platform is free. The content, the tools,
the paid amplification and the staff time are not.\\
\rowcolor{tablealt}
\rlab{``Results are immediate.''} & Community and trust compound slowly. Only
paid reach is instant.\\
\rlab{``You must be on every platform.''} & Presence without competence damages
the brand. Choose platforms by where the persona already is.\\
\rowcolor{tablealt}
\rlab{``More followers means more sales.''} & Vanity metrics do not pay.
Engagement rate and conversion do.\\
\rlab{``It is only for young people, or only for B2C.''} & LinkedIn is the single
most effective B2B channel, and every demographic is now represented.\\
\rowcolor{tablealt}
\rlab{``It is just posting.''} & Listening, responding, moderating and analysing
consume more time than publishing.\\
\rlab{``Anyone can do it.''} & It requires strategy, copywriting, design, data
literacy and crisis judgement.\\
\end{mmlongtable}

\section{Characteristics of a Successful Social Media Marketer}

\begin{itemize}
\item \sfb{Listens more than they broadcast} --- monitors the conversation before
      entering it.
\item \sfb{Understands each platform's native culture} rather than cross-posting
      identical content.
\item \sfb{Is consistent} --- in voice, in visual identity and in publishing
      cadence.
\item \sfb{Is authentic and transparent} --- including, and especially, when
      handling criticism.
\item \sfb{Is data-literate} --- reads analytics and changes behaviour because
      of them.
\item \sfb{Responds fast}, especially to complaints. Social media is a public
      service channel whether or not the brand intends it to be.
\item \sfb{Is creative and adaptive} --- algorithms and formats change
      constantly.
\item \sfb{Builds relationships, not reach} --- with customers, communities and
      influencers.
\end{itemize}

\section{The Customer Lifecycle Orbit}

Social media content only works if it is matched to where the customer already
is. The lifecycle orbit is that map.
\pyq{CIE-II, 26 May 2026 --- 10 M}

\begin{figure}[htbp]
\centering
\begin{tikzpicture}[font=\sffamily\footnotesize]
  \def\R{3.35}
  \foreach \i/\lbl/\ang in {%
      1/{1. AWARENESS}/90,
      2/{2. ENGAGEMENT}/30,
      3/{3. PURCHASE}/-30,
      4/{4. RETENTION\\ \& LOYALTY}/-90,
      5/{5. GROWTH}/-150,
      6/{6. ADVOCACY}/150}{
    \node[draw=ocre,line width=0.9pt,fill=ocrepale,rounded corners=4pt,
          align=center,text width=2.15cm,inner sep=4.5pt,
          font=\sffamily\bfseries\scriptsize,text=slate] (n\i) at (\ang:\R) {\lbl};}
  \foreach \a/\b in {1/2,2/3,3/4,4/5,5/6,6/1}{
    \draw[-{Latex[length=2.3mm]},ocre,line width=1pt]
      (n\a) to[bend left=14] (n\b);}
  \node[circle,draw=greenok,line width=1pt,fill=greenpale,
        align=center,text width=1.95cm,font=\sffamily\bfseries\scriptsize,
        text=greenok,inner sep=2pt] at (0,0) {retention is the profit engine};
\end{tikzpicture}
\caption[The customer lifecycle orbit]%
        {The customer lifecycle orbit. Advocacy returns to awareness, which is
         why the sixth stage is an acquisition channel rather than an endpoint.}
\label{fig:orbit}
\end{figure}

\begin{mmlongtable}{The six stages of the orbit}{tab:orbit}%
{@{}P{0.24\textwidth}P{0.69\textwidth}@{}}
\rowcolor{ocre}\thd{Stage} & \thd{What the brand must do}\\
\midrule
\endfirsthead
\rowcolor{ocre}\thd{Stage} & \thd{What the brand must do}\\
\midrule
\endhead
\rlab{1. Awareness} & Drive brand interest and capture attention\\
\rowcolor{tablealt}
\rlab{2. Engagement} & Listen to behaviours; use mid-stage messages to nudge
towards conversion\\
\rlab{3. Purchase} & Ensure an aligned, frictionless transaction\\
\rowcolor{tablealt}
\rlab{4. Retention and loyalty} & Provide continuous value; resolve issues
rapidly\\
\rlab{5. Growth} & Identify cross-sell and upsell opportunities through targeted
automation\\
\rowcolor{tablealt}
\rlab{6. Advocacy} & Reward loyalty so that customers expand reach through their
own networks\\
\end{mmlongtable}

\begin{keybox}[The Statistics --- Memorise These]
\begin{itemize}
\item Probability of selling to a \textbf{new prospect}: under 20\%.
\item Probability of selling to an \textbf{existing customer}: over 60\%.
\item Repeat customers spend up to \textbf{67\% more}.
\item \textbf{49\%} of marketers achieve higher ROI by focusing on engagement.
\end{itemize}
This is the entire economic argument for retention marketing, and it is why the
orbit has six stages rather than stopping at purchase.
\end{keybox}

\begin{warnbox}[Two Accepted Versions]
The lecture material states the orbit as \emph{Awareness \ra{} Engagement \ra{}
Purchase \ra{} Retention/Loyalty \ra{} Growth \ra{} Advocacy}, with
probabilities of under 20\% against over 60\%.

\smallskip
The released scheme of valuation states it as \emph{Awareness \ra{}
Consideration \ra{} Purchase \ra{} Retention \ra{} Loyalty \ra{} Advocacy}, with
5--20\% for new prospects against 60--70\% for existing ones.

\smallskip
\textbf{Both are accepted.} If unsure, present the six-stage orbit above and
quote the probability as ``roughly 5--20\% for new prospects against 60--70\%
for existing customers'', which satisfies either scheme.
\end{warnbox}

\begin{summarybox}
Social media marketing creates and shares content on social networks to achieve
marketing and branding goals, aiming to convert engagement into traffic, leads,
sales and loyalty. It rests on the trust asymmetry: peers are believed roughly
twice as readily as brands, so the objective is to earn advocacy\ix{advocacy} rather than to
broadcast. Three gears drive business value --- social selling, SEO fuel and
omnichannel synchronisation. Seven myths must be dismantled, chief among them
that social media is free and that follower count is a business metric. The
customer lifecycle orbit runs awareness, engagement, purchase, retention and
loyalty, growth, advocacy --- and because selling to an existing customer is
roughly three times likelier than selling to a new prospect, retention rather
than acquisition is the profit engine.
\end{summarybox}

%%%%%%%%%%%%%%%%%%%%%%%%%%%%%%%%%%%%%%%%%%%%%%%%%%%%%%%%%%%%%%%%%%%%%%%%%%%%%%
\chapter{Major Social Media Platforms for Marketing}
%%%%%%%%%%%%%%%%%%%%%%%%%%%%%%%%%%%%%%%%%%%%%%%%%%%%%%%%%%%%%%%%%%%%%%%%%%%%%%

\syllabusstrip{Major Social Media Platforms for Marketing.}

\begin{learningoutcomes}
\begin{itemize}
\item The platform deployment matrix --- mechanic, best use and key features.
\item The Facebook algorithm filter: what it rewards and what it penalises.
\item Deployment tactics for X, including hashtag practice.
\item The four-layer LinkedIn architecture.
\item Why visual content dominates, with the supporting figures.
\item Six criteria for choosing a platform, and the platform psychology matrix.
\end{itemize}
\end{learningoutcomes}

\section{The Platform Deployment Matrix}

\begin{mmlongtable}{The platform deployment matrix}{tab:platforms}%
{@{}P{0.14\textwidth}P{0.22\textwidth}P{0.30\textwidth}P{0.24\textwidth}@{}}
\rowcolor{ocre}\thd{Platform} & \thd{Core mechanic} & \thd{Best for} & \thd{Key features}\\
\midrule
\endfirsthead
\rowcolor{ocre}\thd{Platform} & \thd{Core mechanic} & \thd{Best for} & \thd{Key features}\\
\midrule
\endhead
\rlab{Facebook} \newline {\scriptsize 2.27 billion users} & User-centric
algorithmic news feed & Building authentic B2C and B2B relationships; highly
targeted paid promotion & Live Video, Groups, Tabs and Apps, Marketplace\\
\rowcolor{tablealt}
\rlab{Instagram} & Visual-first discovery; aspirational identity & Lifestyle,
fashion, food, travel, direct-to-consumer brands; influencer marketing &
Reels, Stories, Shopping tags, Explore\\
\rlab{X (Twitter)} \newline {\scriptsize 330 million users} & Fast-paced
real-time ``virtual watercooler'' & Thought leadership, real-time customer
service, trend participation & Curated Lists, hashtags, live chats, Spaces\\
\rowcolor{tablealt}
\rlab{LinkedIn} \newline {\scriptsize 460 million users} & The world's largest
professional network & Reaching a business mindset, B2B lead generation, talent
recruiting & Showcase Pages, professional groups, long-form publishing\\
\rlab{YouTube} & Dual search-and-social algorithm & Education, product
demonstration, long-form trust building & Playlists, Shorts, chapters, end
screens\\
\rowcolor{tablealt}
\rlab{Pinterest} & Visual search and planning & Home, d\'ecor, DIY, weddings,
recipes --- high purchase intent & Rich Pins, boards\\
\rlab{WhatsApp and Telegram} & Private messaging and broadcast & Customer
support, order updates, community groups & Business catalogue, broadcast lists\\
\end{mmlongtable}

\section{Facebook --- The Algorithm Filter}

\begin{keybox}[The Highest-Probability Ten-Mark Question in This Unit]
The single highest-probability ten-mark question in this unit.
\pyq{CIE-II, 26 May 2026 --- 10 M, both sets}
\end{keybox}

The Facebook algorithm determines which posts appear in users' news feeds. All
published content enters a filter; only some emerges into organic distribution.

\subsection{How Facebook evaluates content}

\begin{enumerate}
\item \sfb{User engagement history} --- what this specific user has previously
      interacted with.
\item \sfb{Relevance} of the content to that user.
\item \sfb{Content type} preferred by that user --- video, image, link or text.
\item \sfb{Timeliness} of the post.
\item \sfb{Quality and originality.}
\end{enumerate}

\begin{figure}[htbp]
\centering
\begin{tikzpicture}[font=\sffamily\footnotesize]
  % published content
  \node[draw=slate,line width=0.9pt,fill=white,rounded corners=4pt,
        minimum width=3.1cm,minimum height=1.0cm,align=center,
        font=\sffamily\bfseries\scriptsize,text=slate] (in) at (0,0)
        {ALL PUBLISHED\\ CONTENT};

  % the filter
  \fill[ocre!22] (2.5,1.5) -- (4.6,0.55) -- (4.6,-0.55) -- (2.5,-1.5) -- cycle;
  \draw[ocre,line width=1pt] (2.5,1.5) -- (4.6,0.55) -- (4.6,-0.55) -- (2.5,-1.5) -- cycle;
  \node[text=ocredark,font=\sffamily\bfseries\scriptsize,align=center]
    at (3.45,0) {THE\\ FILTER};

  \draw[-{Latex[length=2.4mm]},slate,line width=1pt] (in) -- (2.45,0);

  % rewarded
  \node[draw=greenok,line width=0.9pt,fill=greenpale,rounded corners=4pt,
        minimum width=3.6cm,minimum height=1.0cm,align=center,
        font=\sffamily\bfseries\scriptsize,text=greenok] (ok) at (7.3,1.35)
        {PROPELLED THROUGH\\ organic distribution};
  \draw[-{Latex[length=2.4mm]},greenok,line width=1pt] (4.65,0.42) -- (5.5,1.2);

  % penalised
  \node[draw=warnred,line width=0.9pt,fill=warnpale,rounded corners=4pt,
        minimum width=3.6cm,minimum height=1.0cm,align=center,
        font=\sffamily\bfseries\scriptsize,text=warnred] (no) at (7.3,-1.35)
        {FILTERED OUT\\ suppressed reach};
  \draw[-{Latex[length=2.4mm]},warnred,line width=1pt] (4.65,-0.42) -- (5.5,-1.2);

  \node[anchor=north,text=slate,font=\sffamily\scriptsize,align=center,
        text width=11cm] at (4.0,-2.3)
    {evaluated on engagement history \ \textbullet\ \ relevance \ \textbullet\ \
     preferred content type \ \textbullet\ \ timeliness \ \textbullet\ \
     quality and originality};
\end{tikzpicture}
\caption[The Facebook algorithm filter]%
        {The algorithm filter. Publishing is not distribution: content must pass
         five evaluation criteria before any organic reach is granted.}
\label{fig:fbfilter}
\end{figure}

\subsection{What is penalised and what is rewarded}

\begin{mmlongtable}{The Facebook algorithm filter}{tab:fbfilter}%
{@{}P{0.45\textwidth}P{0.48\textwidth}@{}}
\rowcolor{ocre}\thd{Filtered OUT (penalised)} & \thd{Propelled THROUGH (rewarded)}\\
\midrule
\endfirsthead
\rowcolor{ocre}\thd{Filtered OUT (penalised)} & \thd{Propelled THROUGH (rewarded)}\\
\midrule
\endhead
\textbf{Clickbait headlines} --- punished on time-spent-reading ratios and
flagged phrases & \textbf{Time spent} --- seconds spent on a post after it fully
loads, without manipulating engagement through contests\\
\rowcolor{tablealt}
\textbf{Repeated sources} --- the algorithm enforces diversity, reducing
visibility if you post too frequently back-to-back & \textbf{Video actions} ---
unmuting, switching to full-screen or watching in HD count as major success
signals\\
\textbf{Engagement bait} --- ``Like, share, comment now'' & \textbf{Live
broadcasting} --- Facebook Live generates 300\% more watch time than recorded
video\\
\rowcolor{tablealt}
\textbf{Fake news and misinformation} & \textbf{High-quality original content}\\
\textbf{Spammy or repetitive posts} & \textbf{Meaningful interactions} --- real
comments and discussion\\
\rowcolor{tablealt}
\textbf{Low-quality copied content} & \textbf{Informative and engaging posts}\\
\textbf{Excessive promotional posts} & \textbf{Consistent posting}, and content
that generates discussion\\
\end{mmlongtable}

\begin{warnbox}[The Trap in This Topic]
Note that \emph{engagement bait} is penalised while \emph{meaningful
interaction} is rewarded --- even though both produce comments. The algorithm
distinguishes solicited engagement from earned engagement. Asking for a comment
is punished; writing something worth commenting on is rewarded. An answer that
lists ``comments'' as a reward without that distinction misses the point of the
filter.
\end{warnbox}

\section{X (Twitter) --- Deployment Tactics}

\begin{itemize}
\item \sfb{The character constraint.} The limit is 280 characters, but keeping
      posts under 100 characters yields a 17\% higher engagement rate.
\item \sfb{Hashtag arithmetic.} Posts with hashtags get twice the engagement,
      but using more than two causes a 17\% drop. Treat hashtags as targeted
      visual search tags; the optimal count is one or two.
\item \sfb{Link leverage.} Posts containing links are 86\% more likely to be
      reshared. Use a URL shortener to track clicks.
\item \sfb{Strategic curation.} Use Lists, public and private, to segment
      influencers, competitors and partners. This makes curating
      non-promotional content for the 4-1-1 rule (\S\,3.6) highly efficient.
\end{itemize}

\begin{definitionbox}[Hashtag]
A \keyterm{hashtag} is a word or phrase preceded by the \texttt{\#} symbol that
turns text into a clickable, searchable label grouping all posts on that topic.
\pyq{SEE, Oct/Nov 2023 --- 2 M}
\end{definitionbox}

\noindent\textbf{Importance of the hashtag in social media marketing:} it makes
content discoverable beyond your own followers; groups a campaign under one
trackable label; allows participation in trending conversations; enables
user-generated content\ix{user-generated content} collection; supports event and community marketing; and
provides a measurable unit for campaign reporting.

\section{LinkedIn --- The Professional Architecture}

\begin{figure}[htbp]
\centering
\begin{tikzpicture}[font=\sffamily\footnotesize,
    lay/.style={draw=deepblue,line width=0.9pt,minimum width=10.4cm,
                minimum height=1.12cm,align=left,text width=9.9cm,
                font=\sffamily\scriptsize,text=slate,inner sep=6pt}]
  \node[lay,fill=deepbluepale] (l4) at (0,0)      {%
     \textbf{Layer 4 --- Professional Groups (the community).}\ \ Creating or
     joining targeted communities by industry or pain point, to demonstrate
     thought leadership and prospect.};
  \node[lay,fill=white] (l3) at (0,-1.2)          {%
     \textbf{Layer 3 --- Showcase Pages (the niche).}\ \ Dedicated sub-pages
     highlighting specific products or services for distinct audiences.};
  \node[lay,fill=deepbluepale] (l2) at (0,-2.4)   {%
     \textbf{Layer 2 --- Content Publishing (the voice).}\ \ Text posts, slide
     decks and educational content to build followers; educational content
     performs highest.};
  \node[lay,fill=white] (l1) at (0,-3.6)          {%
     \textbf{Layer 1 --- Profile and Company Page (the foundation).}\ \ The
     primary branded landing page; showcases products, the employee network and
     core messaging.};
  \draw[-{Latex[length=2.4mm]},deepblue,line width=1pt]
     (-5.55,-4.05) -- (-5.55,0.45);
  \node[rotate=90,anchor=south,text=deepblue,font=\sffamily\bfseries\scriptsize]
     at (-5.75,-1.8) {increasing depth of relationship};
\end{tikzpicture}
\caption[The four-layer LinkedIn architecture]%
        {LinkedIn's four layers. Each layer presupposes the one beneath it: a
         group presence without a credible company page converts nothing.}
\label{fig:linkedin}
\end{figure}

\section{Visual Content --- Why It Dominates}

\begin{mmlongtable}{The evidence for visual content}{tab:visual}%
{@{}P{0.19\textwidth}P{0.55\textwidth}P{0.18\textwidth}@{}}
\rowcolor{ocre}\thd{Finding} & \thd{Statistic} & \thd{Source}\\
\midrule
\endfirsthead
\rowcolor{ocre}\thd{Finding} & \thd{Statistic} & \thd{Source}\\
\midrule
\endhead
\rlab{Comprehension} & People following directions with text plus illustrations
do 323\% better & NeoMam\\
\rowcolor{tablealt}
\rlab{Virality} & Visual content is 40 times more likely to be shared & Buffer\\
\rlab{Visibility} & Content with relevant images gets 94\% more views &
KISSMetrics\\
\rowcolor{tablealt}
\rlab{Video dominance} & Video generates 135\% greater organic reach than photo
posts & Socialbakers\\
\rlab{Conversion} & 77\% of consumers have been convinced to buy after watching
a video & Wyzowl\\
\end{mmlongtable}

\section{Choosing the Right Platform}

\begin{exambox}[The Comparison Question]
``Compare the roles of Facebook, Instagram, LinkedIn, and Twitter in social
media marketing. How should a business choose the right platform?''
\pyq{CIE-II, 7 May 2025 --- 10 M}
Answer in two movements: the deployment matrix of \S\,2.1 for the comparison,
then the six criteria below for the choice. Close with the psychology matrix.
\end{exambox}

\begin{enumerate}
\item \sfb{Where the persona already is} --- the platform must be one the target
      audience actively uses, not one the marketer likes.
\item \sfb{B2C against B2B} --- Facebook and Instagram skew B2C; LinkedIn is
      decisively B2B.
\item \sfb{Content format capability} --- if the product is visual, Instagram
      and YouTube; if it needs explanation, LinkedIn and YouTube; if it is
      real-time, X.
\item \sfb{Objective and funnel stage} --- awareness on Instagram and YouTube,
      consideration on LinkedIn and YouTube, conversion through Facebook and
      Instagram advertising with retargeting.
\item \sfb{Resources} --- each platform demands a distinct content type. Better
      to dominate two than to be mediocre on six.
\item \sfb{Measurable ROI} --- test with small paid budgets, then concentrate
      spend where cost per result is lowest.
\end{enumerate}

\begin{keybox}[The Platform Psychology Matrix]
Each platform serves a distinct psychological function in the decision process,
and this is the sharpest way to justify a platform choice.

\smallskip
\textbf{TikTok} \ra{} discovery, driven by emotion and novelty. \quad
\textbf{YouTube} \ra{} education, built on retention and perceived expertise.
\quad \textbf{AI assistants} \ra{} verification, demanding citations and
semantic clarity. \quad \textbf{Amazon} \ra{} purchasing, relying on social
proof and trust. \quad \textbf{Instagram} \ra{} aspiration, powered by
aspirational identity. \quad \textbf{Reddit} \ra{} validation, thriving on raw
and unfiltered authenticity. \quad \textbf{Google} \ra{} the click.
\end{keybox}

\begin{summarybox}
Each platform has a core mechanic that determines what it is for: Facebook's
algorithmic feed, Instagram's visual discovery, X's real-time watercooler,
LinkedIn's professional network, YouTube's dual search-and-social algorithm,
Pinterest's visual planning, and messaging apps' private broadcast. The Facebook
filter evaluates engagement history, relevance, preferred content type,
timeliness and originality; it rewards time spent, video actions, Live, original
content and meaningful interaction, and penalises clickbait, repeated sources,
engagement bait, misinformation, spam, copied content and excessive promotion.
On X, brevity, one or two hashtags and links outperform. LinkedIn works in four
layers from company page to groups. Visual and video content dominate every
comparative metric. Choose a platform on persona location, B2C against B2B,
format capability, funnel stage, resources and measured ROI --- and justify it
by the platform's psychological function.
\end{summarybox}

%%%%%%%%%%%%%%%%%%%%%%%%%%%%%%%%%%%%%%%%%%%%%%%%%%%%%%%%%%%%%%%%%%%%%%%%%%%%%%
\chapter{Developing Data-Driven Audience and Campaign Insights}
%%%%%%%%%%%%%%%%%%%%%%%%%%%%%%%%%%%%%%%%%%%%%%%%%%%%%%%%%%%%%%%%%%%%%%%%%%%%%%

\syllabusstrip{Developing Data-driven Audience \& Campaign Insights.}

\begin{learningoutcomes}
\begin{itemize}
\item What data-driven audience insight means, in the words the scheme expects.
\item The eight-step process for developing audience insight.
\item The persona builder canvas.
\item The seven layers by which analytics build an audience profile.
\item The content lifecycle matrix --- objective, tone and format by stage.
\item The 4-1-1 content ratio.
\item The 4~Rs of content scaling, in both taught formulations.
\end{itemize}
\end{learningoutcomes}

\section{What Data-Driven Audience Insight Means}

\begin{exambox}[Model Answer --- 2 Marks]
\pyq{CIE-II, 26 May 2026 --- 2 M}
\keytermas{Data-driven audience insight}{data-driven audience insight} in social
media marketing refers to the process of collecting and analysing user data ---
such as demographics, interests, online behaviour, engagement patterns and
preferences --- from social media platforms in order to understand the target
audience better. It enables marketers to create personalised campaigns and
improve customer engagement.
\end{exambox}

\section{The Eight-Step Insight Process}

\begin{enumerate}
\item \sfb{Define the question.} Decide what decision the insight must support
      --- which platform to invest in, which message to lead with, which segment
      to target. Insight without a decision attached is a report.
\item \sfb{Collect data} from native analytics (Meta and Instagram Insights,
      LinkedIn Page Analytics, X Analytics, YouTube Studio), website analytics,
      CRM records and social listening tools.
\item \sfb{Clean and consolidate} the data into one view, so that
      platform-specific metrics become comparable.
\item \sfb{Segment the audience} by demographics, psychographics, behaviour and
      lifecycle stage.
\item \sfb{Identify patterns} --- peak activity times, best-performing formats,
      highest-engagement topics, common objections and sentiment trends.
\item \sfb{Build audience profiles and personas} from those patterns.
\item \sfb{Apply the insight} --- adjust targeting, creative, tone, format,
      posting schedule and budget allocation.
\item \sfb{Measure and iterate} --- feed campaign results back as new data.
\end{enumerate}

\section{The Persona Builder Canvas}

\begin{mmlongtable}{The persona builder canvas}{tab:canvas2}%
{@{}P{0.26\textwidth}P{0.67\textwidth}@{}}
\rowcolor{ocre}\thd{Quadrant} & \thd{Questions to answer}\\
\midrule
\endfirsthead
\rowcolor{ocre}\thd{Quadrant} & \thd{Questions to answer}\\
\midrule
\endhead
\rlab{Background and goals} & What are their business goals and challenges? What
is their team structure and job role?\\
\rowcolor{tablealt}
\rlab{Information sources} & Where do they consume content? Which social
channels do they actively use --- or avoid?\\
\rlab{Messaging and topics} & Which marketing messages speak to them? Which
content topics capture their interest?\\
\rowcolor{tablealt}
\rlab{Purchase role and objections} & How much influence do they have in the
decision? What objections do they raise during the sales process?\\
\end{mmlongtable}

The working target is four to six distinct personas --- for example Executive
Sponsor, Decision Maker and End User in a B2B setting.

\section{How Social Media Analytics Create Audience Profiles}

\begin{keybox}[High Yield]
This is an eight-mark Semester End Examination question in its own right.
\pyq{SEE, Jun/Jul 2025, Q4a --- 8 M}
\end{keybox}

\begin{figure}[htbp]
\centering
\begin{tikzpicture}[font=\sffamily\footnotesize,
    lyr/.style={draw=ocre,line width=0.8pt,minimum width=10.6cm,
                minimum height=0.86cm,align=left,text width=10.1cm,
                font=\sffamily\scriptsize,text=slate,inner sep=5pt}]
  \node[lyr,fill=ocrepale] (a) at (0,0)     {\textbf{1. Demographic} --- verified age bands, gender split, language, city and country};
  \node[lyr,fill=white]    (b) at (0,-0.92) {\textbf{2. Temporal} --- when followers are online, which sets the publishing schedule};
  \node[lyr,fill=ocrepale] (c) at (0,-1.84) {\textbf{3. Interest} --- page likes, followed accounts and content affinities};
  \node[lyr,fill=white]    (d) at (0,-2.76) {\textbf{4. Behavioural} --- what is saved, shared, commented on, watched to completion};
  \node[lyr,fill=ocrepale] (e) at (0,-3.68) {\textbf{5. Conversational} --- the audience's own vocabulary, complaints and unmet needs};
  \node[lyr,fill=white]    (f) at (0,-4.60) {\textbf{6. Comparative} --- competitor audience overlap, showing whom you fail to reach};
  \node[lyr,fill=greenpale,draw=greenok] (g) at (0,-5.52) {\textbf{7. Conversion} --- social data joined to analytics and CRM: who actually \emph{buys}};
\end{tikzpicture}
\caption[The seven analytic layers of an audience profile]%
        {The seven layers. Layer~4 matters more than layer~3, because saves and
         completions reveal what the audience values rather than what it claims
         to value; layer~7 matters most, because engagement and purchase are not
         the same population.}
\label{fig:sevenlayers}
\end{figure}

\begin{itemize}
\item \sfb{Demographic layer} --- platform analytics supply verified age bands,
      gender split, language, city and country distribution.
\item \sfb{Temporal layer} --- when followers are online determines the
      publishing schedule.
\item \sfb{Interest layer} --- page likes, followed accounts and content
      affinities reveal adjacent interests to target.
\item \sfb{Behavioural layer} --- which posts are saved, shared, commented on or
      watched to completion reveals what the audience actually values, as
      opposed to what it claims to value.
\item \sfb{Conversational layer} --- social listening surfaces the audience's own
      vocabulary, complaints and unmet needs.
\item \sfb{Comparative layer} --- competitor audience overlap shows which
      segments you are failing to reach.
\item \sfb{Conversion layer} --- linking social data to website analytics and
      CRM shows which audience segment actually buys, not merely which one
      engages.
\end{itemize}

\section{The Content Lifecycle Matrix}

\begin{keybox}[Guaranteed Ten-Mark Question]
``Based on the Content Lifecycle Matrix, how should a marketer adjust the
objective, tone, and format of their content as a prospect transitions from the
Early-Stage to the Late-Stage?''
\pyq{CIE-II, 26 May 2026 --- 10 M}
\end{keybox}

\begin{mmlongtable}{The content lifecycle matrix}{tab:lifecyclematrix}%
{@{}P{0.139\textwidth}P{0.257\textwidth}P{0.258\textwidth}P{0.258\textwidth}@{}}
\rowcolor{ocre}\thd{} & \thd{Early-Stage} & \thd{Mid-Stage} & \thd{Late-Stage}\\
\midrule
\endfirsthead
\rowcolor{ocre}\thd{} & \thd{Early-Stage} & \thd{Mid-Stage} & \thd{Late-Stage}\\
\midrule
\endhead
\rlab{Objective} & Nudge the prospect from awareness to engagement & Nudge
towards purchase; build trust & Drive immediate conversions\\
\rowcolor{tablealt}
\rlab{Tone and style} & Fun, entertaining and educational & Educational,
engaging, goal-oriented & Informative, brand-aligned, conversion-focused\\
\rlab{Format examples} & Curated news, helpful tips, visual content & Social
contests, newsletter subscriptions, event invitations & Click-to-purchase, demo
offerings, trial sign-ups\\
\end{mmlongtable}

\begin{figure}[htbp]
\centering
\begin{tikzpicture}[font=\sffamily\footnotesize]
  \draw[-{Latex[length=3mm]},ocre,line width=1.4pt] (0,0) -- (11.2,0);
  \foreach \x/\lbl in {1.4/EARLY,5.6/MID,9.8/LATE}{
    \fill[ocre] (\x,0) circle (0.13);
    \node[anchor=south,text=ocredark,font=\sffamily\bfseries\scriptsize]
      at (\x,0.22) {\lbl-STAGE};}
  \node[anchor=north,text=slate,font=\sffamily\scriptsize,align=center,text width=3.3cm] at (1.4,-0.3)
    {\textit{objective:} attention\\ \textit{tone:} fun\\ \textit{format:} curated, visual};
  \node[anchor=north,text=slate,font=\sffamily\scriptsize,align=center,text width=3.3cm] at (5.6,-0.3)
    {\textit{objective:} trust\\ \textit{tone:} educational\\ \textit{format:} participatory};
  \node[anchor=north,text=slate,font=\sffamily\scriptsize,align=center,text width=3.3cm] at (9.8,-0.3)
    {\textit{objective:} action\\ \textit{tone:} decisive\\ \textit{format:} transactional};
  \node[anchor=north,text=ocredark,font=\sffamily\scriptsize\itshape,align=center,text width=11cm]
    at (5.6,-2.35)
    {broad \ra{} specific \quad\textbullet\quad entertaining \ra{} informative
     \quad\textbullet\quad brand-building \ra{} conversion-driving
     \quad\textbullet\quad ``no ask'' \ra{} ``clear ask''};
\end{tikzpicture}
\caption[The content lifecycle trajectory]%
        {The trajectory the matrix describes. Each of the three rows moves in a
         consistent direction, which is what makes the one-line answer possible.}
\label{fig:lifecycletrajectory}
\end{figure}

\begin{exambox}[The One-Line Answer]
As the prospect moves from early to late stage, content shifts from broad to
specific, from entertaining to informative, from brand-building to
conversion-driving, and from ``no ask'' to ``clear ask''. The \textbf{objective}
moves attention \ra{} trust \ra{} action; the \textbf{tone} moves fun \ra{}
educational \ra{} decisive; the \textbf{format} moves curated and visual \ra{}
interactive and participatory \ra{} transactional.
\end{exambox}

\section{The Content Mix --- The 4-1-1 Rule}

\begin{definitionbox}[The 4-1-1 Rule]
Coined by Joe Pulizzi, the \keytermas{4-1-1 rule}{4-1-1 rule} governs the ratio
of value to promotion in a social feed. For every six posts: \textbf{four}
educational or entertaining posts carrying pure value and no ask, plus
\textbf{one} soft promotion such as a guide or case study, plus \textbf{one}
hard promotion such as a product demonstration or direct offer.
\end{definitionbox}

\begin{figure}[htbp]
\centering
\begin{tikzpicture}[font=\sffamily\footnotesize,
    bx/.style={minimum width=1.55cm,minimum height=1.55cm,align=center,
               font=\sffamily\bfseries\scriptsize,rounded corners=3pt,
               inner sep=2pt},
    sub/.style={anchor=north,align=center,font=\sffamily\scriptsize,
                text width=1.75cm}]

  % Sub-labels are separate nodes rather than nested groups: a \\ inside a
  % braced group within a tikz node breaks the align=center line breaker.
  \foreach \x in {0,1.7,3.4,5.1}{
    \node[bx,draw=greenok,fill=greenpale,text=greenok] at (\x,0) {VALUE};
    \node[sub,text=greenok] at (\x,-0.95) {no ask};}

  \node[bx,draw=ocre,fill=ocrepale,text=ocredark] at (6.8,0) {SOFT\\ PROMO};
  \node[sub,text=ocredark] at (6.8,-0.95) {guide, case study};

  \node[bx,draw=warnred,fill=warnpale,text=warnred] at (8.5,0) {HARD\\ PROMO};
  \node[sub,text=warnred] at (8.5,-0.95) {demo, offer};

  \draw[decorate,decoration={brace,amplitude=5pt},greenok,line width=0.8pt]
    (-0.8,0.92) -- (5.9,0.92);
  \node[anchor=south,text=greenok,font=\sffamily\bfseries\scriptsize] at (2.55,1.1) {4};
  \node[anchor=south,text=ocredark,font=\sffamily\bfseries\scriptsize] at (6.8,0.92) {1};
  \node[anchor=south,text=warnred,font=\sffamily\bfseries\scriptsize] at (8.5,0.92) {1};

  \node[anchor=north,text=slate,font=\sffamily\scriptsize,align=center,text width=11cm]
    at (4.25,-1.75)
    {value far outweighs the pitch: the brand never reads as pushy, and the
     credibility required for late-stage conversion accumulates
     systematically};
\end{tikzpicture}
\caption[The 4-1-1 content ratio]%
        {The 4-1-1 ratio. Two of every six posts may ask for something; four
         may not.}
\label{fig:fouroneone}
\end{figure}

\section{The 4 Rs of Content Scaling}

\begin{keybox}[High Yield]
The 4~Rs let one master asset produce many social assets, which is how a small
team sustains a high posting cadence.
\pyq{CIE-II, 26 May 2026 --- 10 M}
\end{keybox}

\begin{mmlongtable}{The 4 Rs of content scaling}{tab:fourrs}%
{@{}P{0.13\textwidth}P{0.42\textwidth}P{0.38\textwidth}@{}}
\rowcolor{ocre}\thd{R} & \thd{What it means} & \thd{Example}\\
\midrule
\endfirsthead
\rowcolor{ocre}\thd{R} & \thd{What it means} & \thd{Example}\\
\midrule
\endhead
\rlab{Reorganize} & Break a single master asset into smaller micro-pieces &
Extract individual statistics from a large infographic to fuel several separate
social posts\\
\rowcolor{tablealt}
\rlab{Rewrite} & Extend the investment by mining the existing content library;
extract relevant takeaways even from older assets & Turn a two-year-old
whitepaper's key finding into a fresh LinkedIn post\\
\rlab{Retire} & Ruthlessly prune outdated content. Social media is
instantaneous, and outdated information damages brand perception & Delete or
archive posts citing superseded data or discontinued products\\
\rowcolor{tablealt}
\rlab{Redesign} & Refresh the visual packaging; update creative assets to make
existing content relevant for a new segment or persona & Re-skin a B2B case
study's graphics for a student audience\\
\end{mmlongtable}

\begin{keybox}[Why the 4 Rs Matter]
They maximise the return on every content asset: saving time and cost,
increasing reach, improving SEO performance, enhancing audience engagement and
maximising ROI. One master asset can legitimately supply a month of social
content.
\end{keybox}

\begin{warnbox}[Two Accepted Versions]
The lecture material teaches the 4~Rs as \textbf{Reorganize, Rewrite, Retire,
Redesign}.

\smallskip
The released scheme of valuation answers with \textbf{Repurpose, Refresh,
Repost, Recombine} --- repurpose means convert to another format (blog \ra{}
video \ra{} infographic); refresh means update old content with new data; repost
means share high-performing content again at a different time; recombine means
merge several pieces into a new asset such as an e-book or webinar.

\smallskip
\textbf{Safest strategy:} give the lecture version as the primary answer and add
the four alternative terms as ``also described as''. The underlying logic is
identical.
\end{warnbox}

\begin{summarybox}
Data-driven audience insight is the collection and analysis of platform user
data to understand the audience and personalise campaigns. It runs in eight
steps from defining the decision through collection, consolidation,
segmentation, pattern-finding, persona construction and application to
iteration. Analytics build a profile in seven layers, of which the behavioural
and conversion layers matter most. The content lifecycle matrix moves objective
from attention to trust to action, tone from fun to educational to decisive, and
format from curated to participatory to transactional. The 4-1-1 rule caps
promotion at two posts in six. The 4~Rs --- reorganize, rewrite, retire,
redesign, or equivalently repurpose, refresh, repost, recombine --- scale one
master asset into many.
\end{summarybox}

%%%%%%%%%%%%%%%%%%%%%%%%%%%%%%%%%%%%%%%%%%%%%%%%%%%%%%%%%%%%%%%%%%%%%%%%%%%%%%
\chapter{Social Media for Business}
%%%%%%%%%%%%%%%%%%%%%%%%%%%%%%%%%%%%%%%%%%%%%%%%%%%%%%%%%%%%%%%%%%%%%%%%%%%%%%

\syllabusstrip{Social Media for Business.}

\begin{learningoutcomes}
\begin{itemize}
\item The nine business benefits of social media.
\item The ten-step social media marketing plan.
\item The tool categories and what each is for.
\item Blogs, podcasts and webinars compared.
\item Social media monitoring --- metrics, uses and tools.
\end{itemize}
\end{learningoutcomes}

\section{Benefits for Business}

\begin{exambox}[Four Benefits --- 2 Marks]
\pyq{CIE-II, 26 May 2026 --- 2 M}
Increased brand awareness; better customer interaction; cost-effective
promotion; improved customer feedback collection.
\end{exambox}

\begin{itemize}
\item \sfb{Increased brand awareness} and reach beyond the existing customer
      base.
\item \sfb{Better customer interaction} --- a direct, two-way, public channel.
\item \sfb{Cost-effective promotion} relative to traditional media.
\item \sfb{Improved customer feedback collection} --- unsolicited, real-time and
      honest.
\item \sfb{Lead generation and social selling} through personal, relevant
      content.
\item \sfb{Website traffic and indirect SEO benefit} from shares and links.
\item \sfb{Competitive intelligence} --- competitors' campaigns, offers and
      customer complaints are public.
\item \sfb{Crisis detection} --- problems surface on social media before they
      reach support tickets.
\item \sfb{Community and advocacy} --- customers become an unpaid distribution
      network.
\end{itemize}

\section{The Social Media Marketing Plan}

\begin{frameworkbox}[The Ten Steps]
1. Set goals aligned to business objectives \ra{} 2. Audit the current state
\ra{} 3. Define the audience and personas \ra{} 4. Select platforms \ra{}
5. Define the content strategy \ra{} 6. Build the content calendar\ix{content calendar} \ra{}
7. Assign resources and tools \ra{} 8. Publish and engage \ra{} 9. Amplify
\ra{} 10. Measure against KPIs and optimise the next cycle.
\end{frameworkbox}

\begin{enumerate}
\item \sfb{Set goals} aligned to business objectives --- awareness, engagement,
      traffic, leads, sales, retention or service.
\item \sfb{Audit the current state} --- existing profiles, follower quality,
      engagement benchmarks, competitor share of voice\ix{share of voice}.
\item \sfb{Define the audience and personas} (\S\,3.1--3.3).
\item \sfb{Select platforms} using the six criteria of \S\,2.6.
\item \sfb{Define the content strategy} --- themes, formats, tone, the 4-1-1
      mix, and the lifecycle matrix mapping.
\item \sfb{Build the content calendar} --- cadence, timing, campaign moments,
      seasonal hooks.
\item \sfb{Assign resources and tools} --- who creates, who approves, who
      publishes, who responds.
\item \sfb{Publish and engage} --- respond within a defined service-level
      window.
\item \sfb{Amplify} --- paid boosting, employee advocacy, influencer
      partnerships, community seeding.
\item \sfb{Measure against KPIs} and optimise the next cycle.
\end{enumerate}

\section{Social Media Marketing Tools}

\begin{mmlongtable}{Social media marketing tools by category}{tab:smmtools}%
{@{}P{0.24\textwidth}P{0.69\textwidth}@{}}
\rowcolor{ocre}\thd{Category} & \thd{Tools}\\
\midrule
\endfirsthead
\rowcolor{ocre}\thd{Category} & \thd{Tools}\\
\midrule
\endhead
\rlab{Scheduling and management} & Hootsuite, Buffer, Sprout Social, Later, Zoho
Social, Meta Business Suite\\
\rowcolor{tablealt}
\rlab{Design and creative} & Canva, Adobe Express, Figma, CapCut\\
\rlab{Analytics} & Native platform insights, Sprout Social, Socialbakers, Google
Analytics social acquisition reports\\
\rowcolor{tablealt}
\rlab{Listening and monitoring} & Brandwatch, Mention, Talkwalker, Google
Alerts, Hootsuite Streams\\
\rlab{Influencer discovery} & Upfluence, Klear, BuzzSumo\\
\rowcolor{tablealt}
\rlab{Link tracking} & Bitly, UTM parameters via Google's Campaign URL Builder\\
\end{mmlongtable}

\section{Publishing Blogs, Podcasts and Webinars}

\begin{mmlongtable}{Blogs, podcasts and webinars compared}{tab:blogpodwebinar}%
{@{}P{0.13\textwidth}P{0.40\textwidth}P{0.40\textwidth}@{}}
\rowcolor{ocre}\thd{Format} & \thd{Strength} & \thd{Best used for}\\
\midrule
\endfirsthead
\rowcolor{ocre}\thd{Format} & \thd{Strength} & \thd{Best used for}\\
\midrule
\endhead
\rlab{Blog} & Owned, indexable, compounding SEO asset; full control &
Long-tail search capture, thought leadership, evergreen education\\
\rowcolor{tablealt}
\rlab{Podcast} & Intimate; consumed hands-free during commutes and workouts;
builds strong parasocial trust & Deep interviews, category authority, a loyal
niche audience\\
\rlab{Webinar} & Live, interactive, high-intent; captures registration data &
Lead generation, product demonstration, sales enablement, customer training\\
\end{mmlongtable}

\begin{keybox}[Why Webinars Are Strong Marketing Tools]
They generate qualified leads at the point of registration, allow live objection
handling, position the brand as an expert, can be recorded and repurposed
indefinitely through the 4~Rs, and produce measurable attendance and engagement
data.
\end{keybox}

\section{Social Media Monitoring}

\begin{definitionbox}[Social Media Monitoring]
\keytermas{Social media monitoring}{social media monitoring} is the practice of
tracking mentions of your brand, competitors, products, industry keywords and
campaign hashtags across social platforms, in order to respond to individual
conversations and to extract aggregate insight.
\end{definitionbox}

\begin{itemize}
\item \sfb{Key metrics:} mentions, followers, buzz (a combined popularity
      measure), share of voice (conversations about you against conversations
      about the competition) and sentiment (positive, negative or neutral).
\item \sfb{Uses:} reputation management, crisis detection, customer service,
      competitor intelligence, influencer identification and campaign
      measurement.
\item \sfb{Major tools:} Facebook Insights, X Analytics, Hootsuite, Brandwatch,
      Mention, Talkwalker.
\end{itemize}

\begin{summarybox}
Social media delivers nine business benefits, from brand awareness and two-way
interaction to crisis detection and unpaid advocacy. The marketing plan runs ten
steps: goals, audit, audience, platforms, content strategy, calendar, resources,
publish and engage, amplify, measure. Tools divide into scheduling, design,
analytics, listening, influencer discovery and link tracking. Blogs compound as
owned SEO assets, podcasts build parasocial trust, and webinars generate
high-intent leads at registration. Monitoring tracks mentions, followers, buzz,
share of voice and sentiment for reputation management, crisis detection,
service, competitor intelligence and measurement.
\end{summarybox}

%%%%%%%%%%%%%%%%%%%%%%%%%%%%%%%%%%%%%%%%%%%%%%%%%%%%%%%%%%%%%%%%%%%%%%%%%%%%%%
\chapter{Creation and Optimization of Social Media Campaigns}
%%%%%%%%%%%%%%%%%%%%%%%%%%%%%%%%%%%%%%%%%%%%%%%%%%%%%%%%%%%%%%%%%%%%%%%%%%%%%%

\syllabusstrip{Creation \& Optimization of Social Media Campaigns.}

\begin{learningoutcomes}
\begin{itemize}
\item The four elements required for an effective campaign.
\item The ten-step campaign build.
\item Campaign optimisation --- definition, purpose and the eight levers.
\item The six metric families, and why measurement matters.
\item Facebook and Instagram deployed across the funnel.
\end{itemize}
\end{learningoutcomes}

\section{Elements Required for an Effective Campaign}

\begin{exambox}[Model Answer --- 2 Marks]
\pyq{CIE-II, 26 May 2026 --- 2 M}
Four elements required for an effective social media campaign: (1) clear
campaign objectives, (2) a defined target audience, (3) engaging and creative
content, and (4) performance measurement and analytics.
\end{exambox}

\section{Building a Campaign, Step by Step}

\begin{enumerate}
\item \sfb{Set a SMART\ix{SMART} objective} --- for example, ``generate 500 qualified
      leads in eight weeks at under \rupee 200 cost per lead\ix{cost per lead}''.
\item \sfb{Define and segment the target audience} using data-driven insight.
\item \sfb{Choose the platform mix} appropriate to that audience and objective.
\item \sfb{Develop the creative concept} --- a single big idea, a campaign
      hashtag, and a consistent visual identity.
\item \sfb{Produce content mapped to the lifecycle matrix} --- early, mid and
      late-stage assets, scaled through the 4~Rs.
\item \sfb{Build the media plan} --- the organic calendar plus paid budget,
      audiences and bid strategy.
\item \sfb{Set up tracking before launch} --- UTM parameters, conversion pixels,
      goals in analytics. Instrumenting after launch loses the baseline.
\item \sfb{Launch and engage actively} rather than posting and leaving.
\item \sfb{Optimise in flight} (\S\,5.3).
\item \sfb{Report and extract learning} for the next campaign.
\end{enumerate}

\section{Campaign Optimization}

\begin{exambox}[Model Answer --- 2 Marks]
\pyq{CIE-II, 26 May 2026 --- 2 M}
\keytermas{Campaign optimization}{campaign optimization} is the process of
improving social media campaign performance by analysing metrics and adjusting
content, targeting and advertisements.

\smallskip
\textbf{Purpose:} increase engagement; improve conversion rates; reduce
marketing cost; and maximise return on investment\ix{return on investment}.
\end{exambox}

\begin{mmlongtable}{The eight optimisation levers}{tab:levers}%
{@{}P{0.21\textwidth}P{0.72\textwidth}@{}}
\rowcolor{ocre}\thd{Lever} & \thd{Action}\\
\midrule
\endfirsthead
\rowcolor{ocre}\thd{Lever} & \thd{Action}\\
\midrule
\endhead
\rlab{Creative} & A/B test images, video hooks, headlines and captions; retire
fatigued creative\\
\rowcolor{tablealt}
\rlab{Copy and CTA} & Test tone, length and the specific action verb\\
\rlab{Audience and targeting} & Narrow or broaden segments; build lookalike
audiences from converters; exclude existing customers\\
\rowcolor{tablealt}
\rlab{Placement} & Shift budget between feed, Stories, Reels and Explore based
on cost per result\\
\rlab{Timing and frequency} & Publish when the audience is active; cap frequency
to avoid advertising fatigue\\
\rowcolor{tablealt}
\rlab{Budget allocation} & Move spend continuously toward the lowest
cost-per-result ad set\\
\rlab{Landing page} & Fix the destination, not just the advertisement ---
message match, speed and form friction\\
\rowcolor{tablealt}
\rlab{Format} & Move budget toward video, where organic reach and conversion are
strongest\\
\end{mmlongtable}

\begin{warnbox}[Optimise the Destination, Not Only the Advertisement]
The most commonly neglected lever is the landing page. A campaign with excellent
creative and a broken destination reports a targeting problem it does not have.
Always check message match\ix{message match}, load speed and form friction before rebuilding the
audience.
\end{warnbox}

\section{Key Metrics for Campaign Success}

\begin{keybox}[High Yield]
``Discuss the key metrics used to measure the success of a social media
marketing campaign and explain the importance of it for evaluating
performance.''
\pyq{SEE, Oct/Nov 2023, Q3a --- 8 M}
\end{keybox}

\begin{mmlongtable}{The six metric families}{tab:metricfamilies}%
{@{}P{0.19\textwidth}P{0.44\textwidth}P{0.30\textwidth}@{}}
\rowcolor{ocre}\thd{Family} & \thd{Metrics} & \thd{What it tells you}\\
\midrule
\endfirsthead
\rowcolor{ocre}\thd{Family} & \thd{Metrics} & \thd{What it tells you}\\
\midrule
\endhead
\rlab{Reach and volume} & Impressions, reach, follower growth, share of voice &
Whether you are being seen at all\\
\rowcolor{tablealt}
\rlab{Engagement} & Likes, comments, shares, saves, engagement rate, video watch
time, click-through rate & Whether the content resonates\\
\rlab{Traffic} & Sessions from social, pages per session, bounce rate & Whether
attention transfers to owned property\\
\rowcolor{tablealt}
\rlab{Conversion} & Leads, sign-ups, sales, conversion rate, cost per
acquisition, return on ad spend & Whether it produces business value\\
\rlab{Sentiment and advocacy} & Sentiment ratio, brand mentions, user-generated
content volume, net promoter score & Whether perception is improving\\
\rowcolor{tablealt}
\rlab{Efficiency} & Cost per click, cost per thousand impressions, cost per lead
& Whether the money is being spent well\\
\end{mmlongtable}

\begin{exambox}[Why Measurement Matters --- Six Reasons]
It proves ROI and justifies budget; identifies which content and platforms to
scale and which to cut; exposes underperformance early enough to correct it;
enables benchmarking against competitors and against past campaigns; and turns
opinion-driven decisions into evidence-driven ones.
\end{exambox}

\section{Facebook and Instagram Across the Funnel}

\begin{keybox}[High Yield]
``Analyse how businesses can utilize Facebook and Instagram for different stages
of a marketing funnel.''
\pyq{SEE, Jun/Jul 2025, Q3a --- 8 M}
\end{keybox}

\begin{mmlongtable}{Facebook and Instagram across the funnel}{tab:fbigfunnel}%
{@{}P{0.19\textwidth}P{0.37\textwidth}P{0.37\textwidth}@{}}
\rowcolor{ocre}\thd{Funnel stage} & \thd{Facebook} & \thd{Instagram}\\
\midrule
\endfirsthead
\rowcolor{ocre}\thd{Funnel stage} & \thd{Facebook} & \thd{Instagram}\\
\midrule
\endhead
\rlab{Awareness (top)} & Broad-reach video advertising, Page posts, Facebook
Live, Groups seeding & Reels, Explore placement, influencer collaborations,
brand Stories\\
\rowcolor{tablealt}
\rlab{Consideration (middle)} & Lead advertisements, carousel product
advertisements, Messenger conversations, community Groups, event promotion &
Story polls and question stickers, guides, tutorial Reels, saved-post content,
reposted user content\\
\rlab{Conversion (bottom)} & Dynamic product advertisements, retargeting website
visitors and cart abandoners, catalogue advertisements, offer claims &
Instagram Shopping tags, product stickers in Stories, link stickers, direct-message
selling\\
\rowcolor{tablealt}
\rlab{Retention and advocacy (post-purchase)} & Customer Groups, Messenger
support, loyalty offers & Close Friends offers, user-content contests, highlight
reels of customer content\\
\end{mmlongtable}

\begin{summarybox}
An effective campaign needs clear objectives, a defined audience, engaging
creative and measurement. It is built in ten steps, with tracking instrumented
before launch rather than after. Campaign optimisation improves performance by
analysing metrics and adjusting content, targeting and advertisements, working
through eight levers --- creative, copy and CTA, audience, placement, timing,
budget, landing page and format. Measurement runs in six families: reach,
engagement, traffic, conversion, sentiment and efficiency. Facebook and
Instagram each have a distinct role at every funnel stage, from broad-reach video
and Reels at the top to dynamic retargeting and Shopping tags at the bottom.
\end{summarybox}

%%%%%%%%%%%%%%%%%%%%%%%%%%%%%%%%%%%%%%%%%%%%%%%%%%%%%%%%%%%%%%%%%%%%%%%%%%%%%%
%  Unit II, second half: Search Engine Optimization (Chapters 6-11)
%%%%%%%%%%%%%%%%%%%%%%%%%%%%%%%%%%%%%%%%%%%%%%%%%%%%%%%%%%%%%%%%%%%%%%%%%%%%%%

\chapter{Search Engine Optimization Fundamentals}

\syllabusstrip{Search Engine Optimization: Search Engine Optimization Fundamentals.}

\begin{learningoutcomes}
\begin{itemize}
\item What SEO is, and its primary objective.
\item Why SEO matters --- the figures that justify the investment.
\item SEO, PPC and SEM distinguished.
\item The anatomy of a modern search engine results page.
\item The shift from keyword matching to intent, and E-E-A-T.
\item The six-step SEO process, and white-hat against black-hat technique.
\end{itemize}
\end{learningoutcomes}

\section{Definition and Objective}

\begin{definitionbox}[Search Engine Optimization]
\keytermas{Search engine optimization}{search engine optimization} (SEO) is the
set of techniques used to improve the visibility and ranking of a website in the
organic --- that is, unpaid --- results of search engines, thereby increasing the
quantity and quality of traffic to that site.
\end{definitionbox}

\begin{exambox}[Primary Objective --- 2 Marks]
\pyq{SEE, Jun/Jul 2025 --- 2 M}
The primary objective of SEO is to increase organic visibility for relevant
queries so that the right people --- those actively searching with intent ---
find the site, and to increase the site's usability\ix{usability} so that they are satisfied
when they arrive.

\smallskip
\textbf{Visibility without relevance is worthless.} Both halves of that sentence
are required for full marks.
\end{exambox}

\section{Why SEO Matters}

\begin{keybox}[The Case for Organic Search]
\begin{itemize}
\item Around \textbf{10 billion} searches happen every day on Google.
\item \textbf{53\%} of all website traffic is delivered by organic search\ix{organic search} alone.
\item Organic traffic has \textbf{infinite lifespan} --- it compounds, whereas
      paid traffic stops the moment the budget stops.
\item Search traffic carries the \textbf{highest buyer intent} of any channel,
      because the user has explicitly stated their need.
\end{itemize}
\end{keybox}

\noindent\textbf{Importance for online businesses:} lower customer acquisition
cost; credibility, because users trust organic results more than advertisements;
lead generation\ix{lead generation} around the clock; competitive defence; and measurable results.

\section{SEO, SEM and PPC --- The Visibility Matrix}

\begin{mmlongtable}{SEO, PPC and SEM distinguished}{tab:seosempcc}%
{@{}P{0.16\textwidth}P{0.46\textwidth}P{0.31\textwidth}@{}}
\rowcolor{ocre}\thd{Term} & \thd{Mechanism} & \thd{Cost model}\\
\midrule
\endfirsthead
\rowcolor{ocre}\thd{Term} & \thd{Mechanism} & \thd{Cost model}\\
\midrule
\endhead
\rlab{SEO} & Drives organic results and clicks & Free per click; costs time and
resources\\
\rowcolor{tablealt}
\rlab{PPC} \newline {\scriptsize pay per click} & Drives paid results;
advertisers bid on specific keywords & The advertiser is charged per click\\
\rlab{SEM} \newline {\scriptsize search engine marketing} & The umbrella term
combining both organic and paid search strategies & Both\\
\end{mmlongtable}

On a search engine results page, paid listings appear at the top and bottom, with
organic results between them.

\section{The Search Engine Results Page}

\begin{definitionbox}[SERP]
The \keytermas{SERP}{SERP} is the page a search engine returns in response to a
query. It is no longer a simple list of ten blue links.
\end{definitionbox}

\begin{figure}[htbp]
\centering
\begin{tikzpicture}[font=\sffamily\footnotesize,
    el/.style={minimum width=6.4cm,minimum height=0.62cm,align=left,
               text width=6.0cm,font=\sffamily\scriptsize,inner sep=4pt,
               anchor=north west}]
  \draw[rulegrey,line width=0.9pt,rounded corners=3pt] (0,0) rectangle (6.9,-8.5);
  \fill[tablehead,rounded corners=3pt] (0,0) rectangle (6.9,-0.55);
  \node[anchor=west,font=\sffamily\scriptsize,text=slatelight] at (0.18,-0.28)
    {search query};

  \node[el,fill=warnpale,draw=warnred] at (0.25,-0.72)
    {\textbf{\textcolor{warnred}{Paid results}} --- search and shopping ads};
  \node[el,fill=ocrepale,draw=ocre] at (0.25,-1.5)
    {\textbf{\textcolor{ocredark}{Featured snippet}} --- ``position zero''};
  \node[el,fill=deepbluepale,draw=deepblue] at (0.25,-2.28)
    {\textbf{\textcolor{deepblue}{Local pack}} --- three businesses, map, ratings};
  \node[el,fill=white,draw=rulegrey] at (0.25,-3.06)
    {\textbf{Organic result} --- title tag, URL, meta description};
  \node[el,fill=white,draw=rulegrey] at (0.25,-3.84)
    {\textbf{Organic result} --- with rich results: stars, FAQs, price};
  \node[el,fill=greenpale,draw=greenok] at (0.25,-4.62)
    {\textbf{\textcolor{greenok}{People Also Ask}} --- expandable questions};
  \node[el,fill=white,draw=rulegrey] at (0.25,-5.40)
    {\textbf{Image, video and news carousels}};
  \node[el,fill=white,draw=rulegrey] at (0.25,-6.18)
    {\textbf{Organic results} continue};
  \node[el,fill=warnpale,draw=warnred] at (0.25,-6.96)
    {\textbf{\textcolor{warnred}{Paid results}} --- bottom advertisements};
  \node[el,fill=white,draw=rulegrey] at (0.25,-7.74)
    {\textbf{Related searches}};

  \draw[rulegrey,line width=0.9pt,rounded corners=3pt] (7.3,-0.72) rectangle (10.6,-4.3);
  \node[anchor=north,align=center,text width=3.0cm,font=\sffamily\scriptsize,
        text=slate] at (8.95,-0.95)
    {\textbf{Knowledge panel}\\[3pt] entity information drawn from structured
     sources and the knowledge graph};

  \node[anchor=north west,align=left,text width=3.3cm,font=\sffamily\scriptsize,
        text=slatelight] at (7.3,-4.7)
    {\textbf{Three result types}\\[3pt]
     \textcolor{warnred}{paid}\\
     \textcolor{slate}{organic}\\
     \textcolor{deepblue}{universal / vertical}\\ {\scriptsize (images, video,
     news, local, shopping)}};
\end{tikzpicture}
\caption[Anatomy of a modern search engine results page]%
        {A modern SERP. Because paid listings, the featured snippet, the local
         pack and the knowledge panel can all precede the first ordinary organic
         result, ranking first organically no longer means appearing first on the
         page.}
\label{fig:serp}
\end{figure}

\noindent Main elements and result types on a modern SERP:

\begin{itemize}
\item \sfb{Paid results} --- search advertisements at the top and bottom,
      shopping and product listing advertisements, local service advertisements.
\item \sfb{Organic results} --- the ranked list of web pages, each with a title
      tag, a URL and a meta description\ix{meta description}.
\item \sfb{Featured snippet}, known as ``position zero'' --- a direct answer
      extracted from a page.
\item \sfb{Knowledge panel and knowledge graph} --- entity information drawn
      from structured sources.
\item \sfb{Local pack or map pack} --- three local businesses with map, ratings
      and directions.
\item \sfb{People Also Ask} --- expandable related questions, and an excellent
      content-research source.
\item \sfb{Rich results} --- reviews, star ratings, FAQs, recipes and events,
      driven by structured data markup.
\item \sfb{Image, video and news carousels}, and related searches at the foot of
      the page.
\end{itemize}

The three main types of result are therefore \textbf{paid}, \textbf{organic} and
\textbf{universal or vertical} results, the last covering images, video, news,
local and shopping.

\section{The Shift From Keywords to Intent}

Since 2015, Google's \keytermas{RankBrain}{RankBrain} and
\keytermas{Hummingbird}{Hummingbird} algorithms have focused on understanding
user intent rather than matching keyword strings. They expect high-quality
semantic data and well-structured outlines, not keyword stuffing.

\begin{keybox}[E-E-A-T]
Google evaluates content on \textbf{Experience, Expertise, Authoritativeness and
Trustworthiness}. Alongside this runs the link principle: \textbf{link gravity
beats quantity} --- a large quantity of \emph{quality} links is the goal, not
simply a large quantity of links.
\end{keybox}

\begin{warnbox}[The Post-Panda Reality]
Keyword stuffing is dead. Google's sole objective is surfacing the most relevant
answer to a query. If the query is ``how to make banana bread'', the content must
provide the best, most readable, fastest answer.

\smallskip
White-hat optimisation therefore means writing naturally, organising flow well,
and earning backlinks by creating remarkable, link-worthy data or insight.
\end{warnbox}

\section{The Six-Step SEO Process}

An SEO project follows six steps. Note that SEO does not begin and end with
initial work: ongoing success requires monitoring results and building meaningful
content continually.

\begin{figure}[htbp]
\centering
\begin{tikzpicture}[font=\sffamily\footnotesize,
    st/.style={draw=ocre,line width=0.85pt,fill=ocrepale,rounded corners=3pt,
               minimum height=1.5cm,text width=1.72cm,align=center,
               font=\sffamily\bfseries\scriptsize,text=slate,inner sep=3pt}]
  \node[st] (s1) at (0,0)     {1\\ Keyword research};
  \node[st] (s2) at (1.98,0)  {2\\ Reporting \& goal setting};
  \node[st] (s3) at (3.96,0)  {3\\ Content building};
  \node[st] (s4) at (5.94,0)  {4\\ Page optimization};
  \node[st] (s5) at (7.92,0)  {5\\ Social \& link building};
  \node[st] (s6) at (9.90,0)  {6\\ Follow-up reporting};
  \foreach \a/\b in {s1/s2,s2/s3,s3/s4,s4/s5,s5/s6}{
    \draw[-{Latex[length=1.9mm]},ocre,line width=0.85pt] (\a)--(\b);}
  \draw[-{Latex[length=2.1mm]},ocredark,line width=0.85pt,dashed]
    (s6.south) -- ++(0,-0.65) -| (s1.south);
  \node[anchor=north,text=ocredark,font=\sffamily\scriptsize\itshape]
    at (4.95,-1.5) {SEO is a cycle: results are monitored and content built continually};
\end{tikzpicture}
\caption[The six-step SEO process]%
        {The six-step process. Step~6 returns to step~1: an SEO project has no
         completion date, only measurement intervals.}
\label{fig:seoprocess}
\end{figure}

\begin{mmlongtable}{The six steps in detail}{tab:seosteps}%
{@{}P{0.22\textwidth}P{0.71\textwidth}@{}}
\rowcolor{ocre}\thd{Step} & \thd{Activity}\\
\midrule
\endfirsthead
\rowcolor{ocre}\thd{Step} & \thd{Activity}\\
\midrule
\endhead
\rlab{1. Keyword research} & Identify a group of keyword phrases balancing high
usage by searchers with relatively low competition. Also perform competitive
research --- a thorough analysis against the site's seven to ten biggest
competitors using metrics such as indexed content, inbound links, domain age and
social following, to gauge the starting position and identify priorities.\\
\rowcolor{tablealt}
\rlab{2. Reporting and goal setting} & Establish the site's current position in
the search engines and its traffic baseline --- which engines, which phrases,
bounce rates, most popular content. Set measurable goals tied to specific
business objectives.\\
\rlab{3. Content building} & Content is king. High-volume, high-quality content
gives users a reason to stay and return, and gives search engines more
information to store, which translates directly into rankings.\\
\rowcolor{tablealt}
\rlab{4. Page optimization} & On-page work: page titles, text-based navigation,
prominence of targeted keyword phrases, site map, ALT and META data, and cleaning
up the code.\\
\rlab{5. Social and link building} & Establish a social media presence to share
content --- fish where the fish are --- and build high-quality inbound links.
Google views links from high-quality sites as votes for your site.\\
\rowcolor{tablealt}
\rlab{6. Follow-up reporting and analysis} & Repeat the initial reporting at
regular intervals after optimisation. Compare rankings, traffic, social signals
and other key metrics against pre-optimisation levels to produce measurable
results.\\
\end{mmlongtable}

\section{SEO Goals, and White Hat Against Black Hat}

\begin{itemize}
\item The ultimate goal is to \sfb{increase the site's usability} so that the
      right people are brought in from the search engines. The purpose of the
      site must be clearly defined so that you can ensure the site achieves it.
\item When search visitors find that the site meets their expectations after
      clicking through, usability is solid. If they are disappointed, the site
      has missed its mark.
\item \sfb{Follow the best strategies for ranking}, which requires understanding
      how search engines want things done --- familiarise yourself with their
      terms and conditions and webmaster guidelines.
\item \sfb{Use white-hat techniques}, those that fall in line with what search
      engines define as best practice.
\end{itemize}

\begin{warnbox}[Black-Hat Techniques and Their Cost]
Keyword stuffing, cloaking, link farms, hidden text and doorway pages produce
short-term gains and long-term penalties. A manual action or algorithmic penalty
can remove a site from the index entirely, and recovery takes far longer than the
gain ever lasted.
\end{warnbox}

\begin{summarybox}
SEO improves visibility and ranking in organic results to raise the quantity and
quality of traffic; its objective is relevance as well as visibility. It matters
because organic search delivers about 53\% of all web traffic, carries the
highest buyer intent\ix{buyer intent}, compounds indefinitely and costs nothing per click. SEO is
organic, PPC is paid, SEM is the umbrella over both. A modern SERP carries paid
results, organic results and universal results, plus featured snippets,
knowledge panels, local packs, People Also Ask and rich results. Since 2015
RankBrain and Hummingbird have ranked on intent rather than string matching, and
content is judged on E-E-A-T. The process runs six steps --- keyword research,
reporting and goals, content building, page optimisation, social and link
building, follow-up analysis --- and cycles. White-hat technique follows engine
guidelines; black-hat technique trades short gains for long penalties.
\end{summarybox}

%%%%%%%%%%%%%%%%%%%%%%%%%%%%%%%%%%%%%%%%%%%%%%%%%%%%%%%%%%%%%%%%%%%%%%%%%%%%%%
\chapter{Keywords and the SEO Content Plan}
%%%%%%%%%%%%%%%%%%%%%%%%%%%%%%%%%%%%%%%%%%%%%%%%%%%%%%%%%%%%%%%%%%%%%%%%%%%%%%

\syllabusstrip{Keywords and SEO Content Plan.}

\begin{learningoutcomes}
\begin{itemize}
\item Why keyword research matters.
\item Keyword types by intent and by length.
\item The keyword sweet spot.
\item The ten-step SEO content plan, and keyword cannibalisation.
\end{itemize}
\end{learningoutcomes}

\section{Why Keyword Research Matters}

\begin{exambox}[Model Answer --- 2 Marks]
\pyq{CIE-II, 7 May 2025; SEE, Oct/Nov 2023 --- 2 M}
\keytermas{Keyword research}{keyword research} is the process of identifying the
phrases that real users type into search engines, together with their search
volume, competition and intent.

\smallskip
It matters because without it content is optimised for phrases nobody searches,
effort is spent on terms that cannot realistically be ranked for, and traffic
arrives with no intent to buy.
\end{exambox}

\section{Keyword Types}

\subsection{By intent}

\begin{mmlongtable}{Keyword types by search intent}{tab:keywordintent}%
{@{}P{0.24\textwidth}P{0.32\textwidth}P{0.37\textwidth}@{}}
\rowcolor{ocre}\thd{Intent type} & \thd{What the user wants} & \thd{Example}\\
\midrule
\endfirsthead
\rowcolor{ocre}\thd{Intent type} & \thd{What the user wants} & \thd{Example}\\
\midrule
\endhead
\rlab{Informational} & To learn something & ``what is content marketing''\\
\rowcolor{tablealt}
\rlab{Navigational} & To reach a specific site & ``RVCE login''\\
\rlab{Commercial investigation} & To compare before buying & ``best trekking
shoes under 5000''\\
\rowcolor{tablealt}
\rlab{Transactional} & To buy now & ``buy trekking shoes online''\\
\end{mmlongtable}

\subsection{By length}

\begin{figure}[htbp]
\centering
\begin{tikzpicture}[font=\sffamily\footnotesize]
  \fill[ocre!55] (0,0) rectangle (2.1,2.55);
  \node[text=white,font=\sffamily\bfseries\scriptsize,align=center] at (1.05,2.15) {HEAD};
  \node[text=slate,font=\sffamily\scriptsize,align=center,text width=1.95cm] at (1.05,1.15)
    {1--2 words\\ huge volume\\ brutal competition\\ vague intent};

  \fill[ocre!35] (2.3,0) rectangle (4.9,1.75);
  \node[text=slate,font=\sffamily\bfseries\scriptsize,align=center] at (3.6,1.45) {BODY};
  \node[text=slate,font=\sffamily\scriptsize,align=center,text width=2.4cm] at (3.6,0.78)
    {2--3 words\\ moderate volume\\ and competition};

  \fill[ocre!16] (5.1,0) rectangle (11.2,0.95);
  \node[text=slate,font=\sffamily\bfseries\scriptsize,align=center] at (8.15,0.74) {LONG TAIL};
  \node[text=slate,font=\sffamily\scriptsize,align=center,text width=5.9cm] at (8.15,0.3)
    {4+ words \ \textbullet\ \ low volume each but collectively the majority of
     all searches \ \textbullet\ \ low competition, sharp intent, high conversion};

  \draw[-{Latex[length=2.4mm]},slate,line width=0.9pt] (0,-0.4) -- (11.2,-0.4);
  \node[anchor=north,text=slatelight,font=\sffamily\scriptsize] at (5.6,-0.5)
    {increasing phrase length \ \textbullet\ \ decreasing competition \ \textbullet\ \ increasing intent};
\end{tikzpicture}
\caption[Keywords by length: the search demand curve]%
        {The search demand curve. The long tail holds little volume per phrase
         and most of the volume in total, which is why a site with low authority
         competes there rather than at the head.}
\label{fig:longtail}
\end{figure}

\section{Finding the Keyword Sweet Spot}

A keyword is worth targeting only where three conditions overlap.

\begin{figure}[htbp]
\centering
\begin{tikzpicture}[font=\sffamily\footnotesize]
  \begin{scope}[opacity=0.40,transparency group]
    \fill[ocre]     (-1.3,0.6) circle (2.25);
    \fill[deepblue] ( 1.3,0.6) circle (2.25);
    \fill[greenok]  ( 0.0,-1.55) circle (2.25);
  \end{scope}
  \node[text=ocredark,font=\sffamily\bfseries\scriptsize,align=center]
    at (-2.55,1.75) {HIGH SEARCH\\ VOLUME};
  \node[text=deepblue,font=\sffamily\bfseries\scriptsize,align=center]
    at ( 2.55,1.75) {LOW SEO\\ DIFFICULTY};
  \node[text=greenok,font=\sffamily\bfseries\scriptsize,align=center]
    at ( 0.0,-3.4)  {COMMERCIAL INTENT};
  \node[text=slate,font=\sffamily\scriptsize,align=center,text width=2.1cm]
    at (-2.25,0.5) {roughly 250\\ to 1,000+\\ monthly};
  \node[text=slate,font=\sffamily\scriptsize,align=center,text width=2.1cm]
    at ( 2.25,0.5) {under about\\ 40 KD/SD};
  \node[text=slate,font=\sffamily\scriptsize,align=center,text width=2.4cm]
    at ( 0.0,-2.1) {high CPC, buyer-\\ specific phrasing};
  \node[text=white,font=\sffamily\bfseries\scriptsize,align=center]
    at (0,-0.32) {THE SWEET\\ SPOT};
\end{tikzpicture}
\caption[The keyword sweet spot]%
        {Target only the central overlap. Two out of three is a trap: volume
         without intent is vanity, and intent without achievability is wasted
         effort.}
\label{fig:sweetspot2}
\end{figure}

\begin{itemize}
\item \sfb{High search volume} --- a minimum of roughly 250 to 1,000 or more
      monthly searches.
\item \sfb{Low SEO difficulty} --- under about a 40 difficulty score, so that
      ranking is realistically achievable.
\item \sfb{Commercial intent} --- a high cost-per-click and buyer-specific
      phrasing, indicating the searcher intends to transact.
\end{itemize}

\begin{warnbox}[Traffic Without Intent Is Useless]
Do not optimise for clicks; optimise for revenue.
\end{warnbox}

\section{Building an SEO Content Plan}

\begin{enumerate}
\item \sfb{Define business objectives first} --- what does the site actually need
      (sales, leads, sign-ups)?
\item \sfb{Research keywords} and cluster them into topics rather than isolated
      phrases.
\item \sfb{Map intent to funnel stage} --- informational keywords become blog
      content; commercial and transactional keywords become product and
      comparison pages.
\item \sfb{Audit existing content} --- decide what to keep, update, consolidate
      or retire.
\item \sfb{Identify content gaps} against competitors --- topics with demand
      where nobody ranks well.
\item \sfb{Design the site architecture} --- a pillar page per topic cluster,
      with supporting articles internally linked to it.
\item \sfb{Assign one primary keyword per page.}
\item \sfb{Build the editorial calendar} with owners, deadlines and target
      keywords.
\item \sfb{Set measurable targets} --- rankings, impressions, organic sessions,
      assisted conversions.
\item \sfb{Review and iterate} monthly.
\end{enumerate}

\begin{warnbox}[Keyword Cannibalisation]
Never let two of your own pages compete for the same term.
\keytermas{Keyword cannibalisation}{keyword cannibalisation} splits link equity
and click-through between two weaker pages instead of concentrating it in one
strong page, and it leaves the search engine to choose which of your pages to
rank --- a choice it frequently gets wrong. One primary keyword per page,
enforced in the plan, prevents it.
\end{warnbox}

\begin{summarybox}
Keyword research identifies the phrases users actually type, with their volume,
competition and intent; without it, effort goes to phrases nobody searches or
nobody buys from. Intent divides into informational, navigational, commercial
investigation and transactional; length divides into head, body and long tail,
with the long tail holding most total volume at low competition and high
conversion. Target only where volume, achievable difficulty and commercial intent\ix{commercial intent}
overlap. The content plan runs ten steps from business objectives through topic
clustering, intent mapping, content audit, gap analysis, pillar architecture, one
keyword per page, calendar and targets to monthly review.
\end{summarybox}

%%%%%%%%%%%%%%%%%%%%%%%%%%%%%%%%%%%%%%%%%%%%%%%%%%%%%%%%%%%%%%%%%%%%%%%%%%%%%%
\chapter{SEO and Business Objectives}
%%%%%%%%%%%%%%%%%%%%%%%%%%%%%%%%%%%%%%%%%%%%%%%%%%%%%%%%%%%%%%%%%%%%%%%%%%%%%%

\syllabusstrip{SEO \& Business Objectives.}

\begin{learningoutcomes}
\begin{itemize}
\item How to translate each business objective into an SEO focus and a KPI.
\item Why SEO must be designed into website development rather than retro-fitted.
\end{itemize}
\end{learningoutcomes}

\section{Aligning SEO With Business Objectives}

\begin{mmlongtable}{SEO focus and KPI by business objective}{tab:seoobjectives}%
{@{}P{0.22\textwidth}P{0.42\textwidth}P{0.29\textwidth}@{}}
\rowcolor{ocre}\thd{Business objective} & \thd{SEO focus} & \thd{Primary KPI}\\
\midrule
\endfirsthead
\rowcolor{ocre}\thd{Business objective} & \thd{SEO focus} & \thd{Primary KPI}\\
\midrule
\endhead
\rlab{Increase online sales} & Transactional and product keywords; optimised
category and product pages; rich results with reviews and price & Organic
revenue; organic conversion rate\\
\rowcolor{tablealt}
\rlab{Generate leads} & Commercial-investigation keywords; comparison and ``best
X for Y'' content; gated assets & Organic leads; cost per lead\\
\rlab{Build brand authority} & Informational pillar content; E-E-A-T signals;
digital PR and backlinks & Impressions; non-branded keyword coverage; referring
domains\\
\rowcolor{tablealt}
\rlab{Reduce paid advertising dependence} & Rank organically for the terms
currently being bought & Share of traffic from organic against paid\\
\rlab{Enter a new market} & Local and language-specific SEO; localised content
and listings & Local pack visibility; regional organic sessions\\
\end{mmlongtable}

\section{Integrating SEO Into Website Development}

\begin{exambox}[Model Answer --- 2 Marks]
``Why is it important to integrate SEO in website development?''
\pyq{SEE, Jun/Jul 2025 --- 2 M}

\begin{itemize}
\item Site architecture, URL structure and navigation are extremely expensive to
      change after launch but nearly free to get right during the build.
\item Rendering method, page speed and Core Web Vitals\ix{Core Web Vitals} are determined by
      development decisions, not by marketing decisions.
\item Retro-fitting SEO usually forces URL changes, which risk losing
      accumulated ranking authority.
\item Content structure --- heading hierarchy, schema, internal linking --- must
      be designed into the templates.
\item Tracking and analytics must be instrumented before launch, not after
      traffic has been lost.
\item It avoids launching with pages blocked by \texttt{robots.txt} or left
      \texttt{noindex} from staging --- a classic and costly error.
\end{itemize}
\end{exambox}

\begin{summarybox}
Every business objective implies a different SEO focus and a different KPI:
sales imply transactional keywords and organic revenue; leads imply comparison
content and cost per lead\ix{cost per lead}; authority implies pillar content and referring
domains; reducing paid dependence implies ranking for bought terms; new-market
entry implies local and language SEO. SEO belongs in the build, not after it,
because architecture, URLs, rendering, speed, template structure and tracking are
all fixed at development time, and retro-fitting risks the accumulated authority
of existing URLs.
\end{summarybox}

%%%%%%%%%%%%%%%%%%%%%%%%%%%%%%%%%%%%%%%%%%%%%%%%%%%%%%%%%%%%%%%%%%%%%%%%%%%%%%
\chapter{Writing SEO Content}
%%%%%%%%%%%%%%%%%%%%%%%%%%%%%%%%%%%%%%%%%%%%%%%%%%%%%%%%%%%%%%%%%%%%%%%%%%%%%%

\syllabusstrip{Writing SEO Content.}

\begin{learningoutcomes}
\begin{itemize}
\item The anatomy of a top-ranking page.
\item The nine-point method for integrating keywords naturally.
\item Seven considerations when writing SEO-friendly content.
\end{itemize}
\end{learningoutcomes}

\section{Anatomy of a Top-Ranking Page}

\begin{itemize}
\item \sfb{H1 title} --- click-worthy, with the exact-match keyword.
\item \sfb{Introduction} --- comprehensive, answering the user's core intent
      immediately rather than after four hundred words of preamble.
\item \sfb{Structure} --- H2 and H3 subheadings containing semantically related
      long-tail keyword\ix{long-tail keyword} variants.
\item \sfb{Formatting} --- short, scannable paragraphs and bulleted lists.
\item \sfb{Rich media} --- embedded videos and infographics to increase dwell
      time.
\item \sfb{Architecture} --- internal links to keep users inside your ecosystem,
      plus credible outbound links to build trust.
\end{itemize}

\section{Integrating Keywords Naturally}

\begin{exambox}[The Scenario Question]
``You are writing a blog post about a complex technical topic. Describe how you
would integrate relevant keywords naturally into the content while ensuring it
remains engaging and informative for readers.''
\pyq{CIE-II, 24 Jul 2024 --- 10 M}
\end{exambox}

\begin{enumerate}
\item \sfb{Write for the reader first, then optimise.} Draft the piece to
      genuinely answer the question, then place keywords during editing.
\item \sfb{Place the primary keyword in the high-value positions} --- the title
      tag, the H1, the first hundred words, one H2, the URL slug, the meta
      description and one image ALT attribute.
\item \sfb{Use semantic variants and synonyms} rather than repeating the exact
      phrase; modern algorithms understand meaning, not string matching.
\item \sfb{Respect keyword density.} Industry practice is roughly 3--5\% for the
      focus phrase --- enough to signal the topic, not enough to read as spam.
\item \sfb{Answer the sub-questions surfaced in ``People Also Ask''}; each
      becomes an H2 or H3 that naturally carries a long-tail variant.
\item \sfb{Explain jargon on first use}, and use analogies, examples, diagrams
      and code snippets to keep a technical topic engaging.
\item \sfb{Break up the text} with subheadings, lists, tables and visuals every
      few paragraphs.
\item \sfb{Link internally} with keyword-rich but descriptive anchor text\ix{anchor text}.
\item \sfb{Read it aloud.} If a sentence sounds unnatural because of a keyword,
      rewrite it --- user signals now outweigh keyword placement.
\end{enumerate}

\section{Key Considerations}

\begin{exambox}[Two Considerations --- 2 Marks]
\pyq{CIE-II Quiz, 24 Jul 2024 --- 2 M}
(1) \textbf{Match the search intent} --- the format and depth must answer what
the user actually wants. (2) \textbf{Integrate keywords naturally} in the title,
H1, opening paragraph and subheadings using semantic variants, at roughly 3--5\%
density, without stuffing. Also acceptable: a unique title tag and meta
description; a scannable structure.
\end{exambox}

\begin{itemize}
\item \sfb{Search intent match} --- the format must match what already ranks: a
      list, a guide, a comparison, a tool.
\item \sfb{Originality and depth} --- content must be \emph{better} than what
      currently occupies the SERP, not merely equivalent.
\item \sfb{Readability} --- short sentences, plain language, scannable structure.
\item \sfb{Freshness} --- update statistics and examples; stale content loses
      rankings.
\item \sfb{Unique title tags and meta descriptions} on every page.
\item \sfb{Accessibility} --- descriptive ALT text on every meaningful image.
\item \sfb{Comprehensiveness} --- cover the topic fully, so the user has no
      reason to return to the SERP.
\end{itemize}

\begin{summarybox}
A top-ranking page has a keyword-bearing H1, an introduction that answers the
intent immediately, H2 and H3 structure carrying semantic variants, scannable
formatting, rich media for dwell time\ix{dwell time}, and both internal and credible outbound
links. Integrate keywords by writing for the reader first and optimising second,
placing the primary phrase in the high-value positions, using semantic variants
at roughly 3--5\% density, answering People Also Ask questions as subheadings,
explaining jargon, breaking up text, linking internally with descriptive anchors,
and reading aloud to catch unnatural phrasing. Content must match intent, exceed
what already ranks, read easily, stay fresh, carry unique metadata, be accessible
and be complete.
\end{summarybox}

%%%%%%%%%%%%%%%%%%%%%%%%%%%%%%%%%%%%%%%%%%%%%%%%%%%%%%%%%%%%%%%%%%%%%%%%%%%%%%
\chapter{On-Site and Off-Site SEO}
%%%%%%%%%%%%%%%%%%%%%%%%%%%%%%%%%%%%%%%%%%%%%%%%%%%%%%%%%%%%%%%%%%%%%%%%%%%%%%

\syllabusstrip{On-site \& off-site SEO.}

\begin{learningoutcomes}
\begin{itemize}
\item On-page SEO and its five key elements.
\item Nine interlinking best practices.
\item Technical SEO --- the defensive foundation.
\item Off-page SEO, backlinks and brand signals.
\item The on-site against off-site comparison.
\item The House of SEO --- how the three layers combine.
\end{itemize}
\end{learningoutcomes}

\section{On-Site (On-Page) SEO}

\begin{definitionbox}[On-Page SEO]
\keytermas{On-page SEO}{on-page SEO} refers to how well a website's own content
and code are presented to search engines. It involves ensuring that a particular
webpage is structured so that it is found for given keywords and key phrases. It
improves search ranking and the overall readability of the site, and --- unlike
off-page work --- it can be improved immediately by fixing incorrect elements on
the page.
\end{definitionbox}

\begin{mmlongtable}{The five key elements of on-page SEO}{tab:onpage}%
{@{}P{0.16\textwidth}P{0.77\textwidth}@{}}
\rowcolor{ocre}\thd{Element} & \thd{What to do}\\
\midrule
\endfirsthead
\rowcolor{ocre}\thd{Element} & \thd{What to do}\\
\midrule
\endhead
\rlab{1. Page copy} & Produce original, unique, high-quality, relevant content
continuously. Focus each piece on a single primary keyword or key phrase,
maintaining a keyword density of about 3--5\%, mixing primary and secondary
phrases. Quality over quantity --- depth and content intensity have outperformed
sheer length since the 2011 algorithm update.\\
\rowcolor{tablealt}
\rlab{2. Title tags} & Arguably the most important element of the ``big three'',
the others being page copy and inbound links. The title tag supplies the
clickable link in the search result. Google limits page titles to about 70
characters, so titles must be keyword-relevant yet concise. Order: primary
keyword, then secondary keyword, then branded keyword at the end --- except on
the home page.\\
\rlab{3. Meta data} & A well-written description summarising the page, limited to
155--160 characters. Meta keywords lost ranking influence after Google's
September 2009 declaration because of spam abuse, though crawlers still read them
for topicality. The meta description \emph{cannot} influence rankings but
strongly influences click-through rate --- it works as advertising copy for the
organic result, and Google highlights matching keywords in bold.\\
\rowcolor{tablealt}
\rlab{4. Heading tags} & Define the content section by section, like traditional
headings and subheadings. There should be one \texttt{<h1>} per page containing
the most relevant key phrase. Tags run to \texttt{<h6>}, though general practice
uses up to \texttt{<h3>}. Place other important key phrases in \texttt{<h2>} and
\texttt{<h3>}.\\
\rlab{5. Interlinking} & Linking one page of your site to related pages provides
context to both search engines and readers.\\
\end{mmlongtable}

\subsection{Interlinking best practices}

\begin{enumerate}
\item Include links in the main content of each page.
\item Paragraph links carry the most weight.
\item Use keyword-rich anchor text.
\item Avoid non-descriptive anchor text such as ``read more'' or ``click here''.
\item Link to relevant, deep pages --- not only the home page.
\item Use breadcrumb navigation on every page.
\item Monitor inbound links through Google Search Console\ix{Google Search Console}.
\item Avoid multiple links to the same page from a single page.
\item Fewer links means more authority per link.
\end{enumerate}

\section{Technical SEO --- The Defence}

\begin{quotebox}
SEO is like a sports team. To win, both defence and offence are required.
\end{quotebox}

Technical SEO is the foundation of the house: the architecture that allows the
site to be crawled and indexed at all.

\begin{itemize}
\item Site speed and \sfb{Core Web Vitals} --- LCP, INP and CLS.
\item Mobile friendliness and responsive design\ix{responsive design}.
\item HTTPS security.
\item An XML sitemap and a correct \texttt{robots.txt}.
\item Crawlability and indexing\ix{indexing} --- no orphan pages, no broken links, no
      redirect chains.
\item Clean URL structure --- short, readable, keyword-bearing.
\item Canonical tags, to prevent duplicate-content dilution.
\item Structured data and schema markup, for rich results.
\item Avoiding intrusive interstitials and pop-ups that block content on mobile.
\end{itemize}

\begin{exambox}[Two Technical Aspects --- 2 Marks]
``Name two technical aspects of a website that can significantly impact its
organic search ranking.''
\pyq{CIE-II Quiz, 24 Jul 2024 --- 2 M}
Best two answers: \textbf{site loading speed} (Core Web Vitals) and
\textbf{mobile responsiveness}. Also acceptable: HTTPS security, crawlability and
indexability, XML sitemap, URL structure.
\end{exambox}

\section{Off-Site (Off-Page) SEO}

\begin{definitionbox}[Off-Page SEO]
\keytermas{Off-page SEO}{off-page SEO} refers to a website's overall authority on
the web, determined by what other websites say about it. It is a long-term
process. Put simply, off-page is about your online reputation. Its central
activity is acquiring \keytermas{backlinks}{backlink} from authority sites in
your niche --- backlinks are the currency of any off-page strategy. Unlike
on-page work, off-page effort is not visible on the webpage itself; it does the
background work for a better search result.
\end{definitionbox}

\subsection{The two halves of off-page work}

\noindent\sfb{1. Acquiring backlinks.} Search engines treat link popularity as a
key ranking factor, though it is no longer the top factor, because it can be
manipulated. After the Google Panda and Penguin updates, many effective
old-school practices became obsolete and negatively affected large, high-ranking
websites. Engines now focus more on content quality and engagement level than on
the raw number of links. The success factor is not building a long list of
inbound links but building a trail of quality links.

\begin{keybox}[The Link Rule]
Getting links from multiple \emph{relevant} domains in your industry is the key.
\textbf{Link gravity beats quantity} --- a large quantity of quality links is the
goal.
\end{keybox}

\noindent\sfb{2. Building brand signals} --- brand mentions, reviews, ratings,
social sharing, PR coverage and demand generation, all of which contribute to
E-E-A-T.

\subsection{Off-page techniques}

\begin{itemize}
\item \sfb{Digital PR and link earning} --- publish original research, data or
      tools that journalists and bloggers cite naturally.
\item \sfb{Guest posting} on relevant, authoritative sites in the niche.
\item \sfb{Broken-link building} --- find dead links on relevant pages and offer
      your resource as a replacement.
\item \sfb{Business listings and citations} --- Google Business Profile, industry
      directories, consistent name, address and phone details.
\item \sfb{Online reviews and ratings} on Google, industry platforms and
      marketplaces.
\item \sfb{Social media sharing} --- not a direct ranking factor, but it produces
      the exposure from which links follow.
\item \sfb{Influencer and expert collaborations} --- interviews, quotes and
      co-created content.
\item \sfb{Forum and community participation} --- contributing genuine value in
      niche communities.
\item \sfb{Unlinked-mention reclamation} --- find places where the brand is
      mentioned without a link, and request one.
\end{itemize}

\section{On-Site Against Off-Site SEO}

\begin{keybox}[High Yield]
Asked both as a two-mark distinction and as a ten-mark analysis in the same
paper.
\pyq{CIE-II, 7 May 2025 --- 2 M and 10 M}
\end{keybox}

\begin{mmlongtable}{On-site against off-site SEO}{tab:onvsoff}%
{@{}P{0.15\textwidth}P{0.39\textwidth}P{0.39\textwidth}@{}}
\rowcolor{ocre}\thd{Basis} & \thd{On-Site (On-Page) SEO} & \thd{Off-Site (Off-Page) SEO}\\
\midrule
\endfirsthead
\rowcolor{ocre}\thd{Basis} & \thd{On-Site (On-Page) SEO} & \thd{Off-Site (Off-Page) SEO}\\
\midrule
\endhead
\rlab{Location of work} & Within your own website & Outside your website\\
\rowcolor{tablealt}
\rlab{What it controls} & How content and code are presented to search engines &
The site's authority and reputation across the web\\
\rlab{Primary target} & Humans and crawlers reading the page & Third-party
validators --- other sites and users\\
\rowcolor{tablealt}
\rlab{Core challenge} & Quality and formatting & Brand authority\\
\rlab{Degree of control} & Complete --- you own it & Limited --- others decide\\
\rowcolor{tablealt}
\rlab{Speed of effect} & Immediate; can be fixed today & Slow; a long-term
process\\
\rlab{Techniques} & Title tags, meta descriptions, heading tags, keyword
placement, content quality, internal linking, image ALT text, URL structure,
site speed, mobile friendliness & Backlink building, guest posting, digital PR,
social sharing, influencer collaboration, online reviews, business listings,
forum participation\\
\rowcolor{tablealt}
\rlab{Example} & Rewriting a product page title to ``Trekking Shoes for Men ---
Waterproof --- BrandName'' & Earning a link from a national outdoors magazine's
``best trekking gear'' article\\
\end{mmlongtable}

\section{The House of SEO}

\begin{figure}[htbp]
\centering
% Each line of each storey is a separately anchored node at a fixed y, so that
% text can never grow into the storey below it.
\begin{tikzpicture}[font=\sffamily\footnotesize,
    ln/.style={anchor=north,align=center,font=\sffamily\scriptsize,text=slate}]

  % ---- roof: off-page SEO
  \fill[greenpale]              (0,4.3) -- (5.6,6.1) -- (11.2,4.3) -- cycle;
  \draw[greenok,line width=1pt] (0,4.3) -- (5.6,6.1) -- (11.2,4.3) -- cycle;
  \node[ln,text=greenok,font=\sffamily\bfseries\small] at (5.6,5.72) {OFF-PAGE SEO};
  \node[ln,font=\sffamily\scriptsize\itshape]           at (5.6,5.26) {the roof --- the offence};
  \node[ln,text width=8.0cm]                            at (5.6,4.90)
    {link building, PR, ratings, demand generation, E-E-A-T};

  % ---- living floor: on-page SEO
  \fill[ocrepale]            (0.55,1.35) rectangle (10.65,4.3);
  \draw[ocre,line width=1pt] (0.55,1.35) rectangle (10.65,4.3);
  \node[ln,text=ocredark,font=\sffamily\bfseries\small] at (5.6,4.18) {ON-PAGE SEO};
  \node[ln,font=\sffamily\scriptsize\itshape]           at (5.6,3.74) {the living floor};
  \node[ln,text width=9.4cm]                            at (5.6,3.38)
    {relevant topics, title tags, H1--H6 headings, image ALT text, multimedia};
  \node[ln,text width=9.4cm]                            at (5.6,2.72)
    {\textbf{goal:} publish original, error-free content better than the SERP
     competitors};

  % ---- foundation: technical SEO
  \fill[deepbluepale]            (0.55,-1.15) rectangle (10.65,1.35);
  \draw[deepblue,line width=1pt] (0.55,-1.15) rectangle (10.65,1.35);
  \node[ln,text=deepblue,font=\sffamily\bfseries\small] at (5.6,1.23) {TECHNICAL SEO};
  \node[ln,font=\sffamily\scriptsize\itshape]           at (5.6,0.79) {the foundation --- the defence};
  \node[ln,text width=9.4cm]                            at (5.6,0.43)
    {URL structure, HTTPS, mobile friendliness, Core Web Vitals, fast hosting};
  \node[ln,text width=9.4cm]                            at (5.6,-0.23)
    {\textbf{goal:} complete crawlability, indexing and seamless UX
     infrastructure};

  % ---- reading of the diagram
  \node[ln,text=slatelight,font=\sffamily\scriptsize\itshape,text width=11cm]
    at (5.6,-1.45)
    {technical makes the site discoverable, on-page makes it relevant, off-page
     makes it authoritative --- and rankings are a function of all three
     together};
\end{tikzpicture}
\caption[The House of SEO]%
        {The House of SEO. The layers are not alternatives: authority cannot
         compensate for a site that cannot be crawled, and crawlability earns
         nothing without content worth ranking.}
\label{fig:houseofseo}
\end{figure}

\begin{mmlongtable}{The three layers compared}{tab:threelayers}%
{@{}P{0.147\textwidth}P{0.255\textwidth}P{0.255\textwidth}P{0.255\textwidth}@{}}
\rowcolor{ocre}\thd{} & \thd{Technical SEO} & \thd{On-Page SEO} & \thd{Off-Page SEO}\\
\midrule
\endfirsthead
\rowcolor{ocre}\thd{} & \thd{Technical SEO} & \thd{On-Page SEO} & \thd{Off-Page SEO}\\
\midrule
\endhead
\rlab{Position} & The foundation & The living floor & The roof\\
\rowcolor{tablealt}
\rlab{Primary target} & Search engine bots & Humans and algorithms &
Third-party validators\\
\rlab{Core challenge} & Discoverability and security & Quality and formatting &
Brand authority\\
\rowcolor{tablealt}
\rlab{The goal} & Complete crawlability, indexing and seamless UX
infrastructure & Publish well-written, original, error-free content better than
SERP competitors & Enhance brand awareness and recognition; acquire high-quality
backlinks\\
\rlab{Elements} & URL structure, HTTPS, mobile friendliness, Core Web Vitals,
fast hosting, avoiding intrusive interstitials & Relevant topics, title tags,
H1--H6 headings, image ALT text, multimedia & Link building, PR, ratings, demand
generation, E-E-A-T\\
\end{mmlongtable}

\begin{summarybox}
On-page SEO governs how a site's own content and code are presented, through five
elements --- page copy at 3--5\% keyword density, title tags under about 70
characters, meta descriptions of 155--160 characters that drive click-through
rather than ranking, one H1 per page with H2 and H3 beneath, and interlinking with
descriptive anchor text. Technical SEO is the defensive foundation: speed, Core
Web Vitals, mobile friendliness, HTTPS, sitemap, crawlability, clean URLs,
canonicals and schema. Off-page SEO is authority earned elsewhere, principally
through quality backlinks --- link gravity beats quantity --- plus brand signals.
On-page is controllable and immediate; off-page is uncontrollable and slow. The
House of SEO combines all three: technical makes a site discoverable, on-page
makes it relevant, off-page makes it authoritative.
\end{summarybox}

%%%%%%%%%%%%%%%%%%%%%%%%%%%%%%%%%%%%%%%%%%%%%%%%%%%%%%%%%%%%%%%%%%%%%%%%%%%%%%
\chapter{Optimizing Organic Search Ranking}
%%%%%%%%%%%%%%%%%%%%%%%%%%%%%%%%%%%%%%%%%%%%%%%%%%%%%%%%%%%%%%%%%%%%%%%%%%%%%%

\syllabusstrip{Optimize Organic Search Ranking.}

\begin{learningoutcomes}
\begin{itemize}
\item The purpose of optimising organic rankings.
\item Eleven metrics by which organic ranking is measured.
\item The nine-step SEO audit, and a worked diagnostic example.
\item Why good design is now good SEO.
\item The SEO toolset.
\end{itemize}
\end{learningoutcomes}

\section{Purpose}

\begin{exambox}[Model Answer --- 2 Marks]
\pyq{CIE-II, 7 May 2025 --- 2 M}
The purpose of optimising organic search rankings is to appear higher in unpaid
search results for queries with real commercial intent, thereby capturing
sustainable, high-quality, near-zero-marginal-cost traffic; reducing dependence
on paid advertising; and building the credibility that comes with ranking
organically rather than paying for position.
\end{exambox}

\section{Key Metrics}

\begin{mmlongtable}{Metrics for organic search performance}{tab:organicmetrics}%
{@{}P{0.30\textwidth}P{0.63\textwidth}@{}}
\rowcolor{ocre}\thd{Metric} & \thd{What it measures}\\
\midrule
\endfirsthead
\rowcolor{ocre}\thd{Metric} & \thd{What it measures}\\
\midrule
\endhead
\rlab{Average position and keyword rankings} & Where the site sits for its
target keywords\\
\rowcolor{tablealt}
\rlab{Organic impressions} & How often pages are shown in search results\\
\rlab{Organic clicks and organic CTR} & How often people actually click --- and
how compelling the title and meta description are\\
\rowcolor{tablealt}
\rlab{Organic sessions} & Volume of visitors arriving from unpaid search\\
\rlab{Keyword coverage} & The number of distinct keywords ranking in the top 3,
top 10 and top 100\\
\rowcolor{tablealt}
\rlab{Indexed pages} & How much of the site the engine has actually stored\\
\rlab{Referring domains and backlink quality} & Off-page authority growth\\
\rowcolor{tablealt}
\rlab{Domain and page authority} & Third-party authority scoring\\
\rlab{Bounce rate and dwell time} & Whether the traffic finds what it expected\\
\rowcolor{tablealt}
\rlab{Organic and assisted conversions} & Whether the traffic produces business
value\\
\rlab{Core Web Vitals} & Technical experience signals feeding into ranking\\
\end{mmlongtable}

\section{Conducting an SEO Audit}

\begin{keybox}[High Yield]
Examined together with the metrics of \S\,11.2.
\pyq{CIE-II, 24 Jul 2024 --- 10 M}
\end{keybox}

\begin{enumerate}
\item \sfb{Crawl the site} --- with Screaming Frog\ix{Screaming Frog}, SEMrush Site Audit or Ahrefs
      Site Audit --- to find broken links, redirect chains, duplicate titles and
      descriptions, missing H1s, thin pages and orphan pages.
\item \sfb{Check indexation} in Google Search Console. Are the important pages
      actually indexed? Are unimportant ones consuming crawl\ix{crawl} budget?
\item \sfb{Audit technical health} --- speed, Core Web Vitals, mobile usability,
      HTTPS, sitemap, \texttt{robots.txt}, canonical tags, structured data.
\item \sfb{Audit on-page elements} --- title tags, meta descriptions, heading
      hierarchy, keyword targeting, keyword cannibalisation, image ALT text,
      internal linking.
\item \sfb{Audit content} --- relevance, depth, freshness, originality, intent
      match, and gaps against competitors.
\item \sfb{Audit backlinks} --- referring domains, anchor-text distribution,
      toxic links to disavow, and the link gap against competitors.
\item \sfb{Benchmark competitors} on the same metrics, to establish relative
      position.
\item \sfb{Prioritise findings by impact against effort}, and produce a ranked
      action plan.
\item \sfb{Implement, then re-measure} at defined intervals against the
      pre-optimisation baseline.
\end{enumerate}

\begin{examplebox}[Worked Diagnostic: an Institution That Cannot Rank for ``Online Courses'']
\pyq{CIE-II, 24 Jul 2024 --- 10 M}

\textbf{Likely issues.} Targeting head keywords far above the site's authority;
thin course pages with duplicate descriptions; no topic-cluster architecture;
slow, unoptimised pages; no schema markup for courses; very few referring domains
compared with established competitors; and no student-facing informational
content capturing early-funnel search.

\textbf{Recommended actions.}
\begin{enumerate}
\item Shift targeting to achievable long-tail phrases --- ``online data
      analytics course for working professionals in Bengaluru''.
\item Build a pillar page per subject area with supporting articles internally
      linked.
\item Rewrite every course page with unique titles, meta descriptions, outcomes,
      fees, faculty and FAQs.
\item Add Course and FAQ schema for rich results.
\item Fix speed and mobile usability.
\item Earn links through original placement-outcome research, faculty guest
      articles and local news coverage.
\item Build reviews on Google Business Profile and education directories.
\item Measure keyword coverage, organic sessions and enquiry conversions
      monthly.
\end{enumerate}
\end{examplebox}

\section{The Human--Algorithm Symbiosis}

\begin{keybox}[Why Good Design Is Now Good SEO]
Algorithms care about crawlability, keywords, metadata and links. Humans care
about scanning, speed, usability and satisfaction. The \textbf{overlap} is Core
Web Vitals --- speed and stability --- and E-E-A-T --- trust and authority.

\smallskip
Search engines have evolved to mimic human preferences. Good web usability is now
a fundamental requirement for good SEO, not a separate concern.
\end{keybox}

\section{SEO Tools}

\begin{mmlongtable}{SEO tools by category}{tab:seotools}%
{@{}P{0.24\textwidth}P{0.69\textwidth}@{}}
\rowcolor{ocre}\thd{Category} & \thd{Tools}\\
\midrule
\endfirsthead
\rowcolor{ocre}\thd{Category} & \thd{Tools}\\
\midrule
\endhead
\rlab{All-in-one suites} & SEMrush, Ahrefs, Moz Pro, Ubersuggest, SE Ranking\\
\rowcolor{tablealt}
\rlab{Google's own tools} & Google Search Console, Google Analytics 4, Keyword
Planner, PageSpeed Insights, Trends, Rich Results Test, Tag Manager\\
\rlab{Keyword research} & Keyword Planner, SEMrush Keyword Magic Tool, Ahrefs
Keywords Explorer, AnswerThePublic, AlsoAsked\\
\rowcolor{tablealt}
\rlab{Technical crawling} & Screaming Frog SEO Spider, Sitebulb, DeepCrawl\\
\rlab{Backlink analysis} & Ahrefs, Majestic, Moz Link Explorer, SEMrush Backlink
Audit\\
\rowcolor{tablealt}
\rlab{Rank tracking} & SEMrush Position Tracking, AccuRanker, SERPWatcher\\
\rlab{Content optimisation} & Surfer SEO, Clearscope, Yoast SEO, Rank Math\\
\rowcolor{tablealt}
\rlab{Competitive intelligence} & SimilarWeb, SpyFu, BuzzSumo\\
\end{mmlongtable}

\begin{summarybox}
Optimising organic rankings captures sustainable high-intent traffic at near-zero
marginal cost, reduces paid dependence and confers the credibility of an organic
listing. Performance is measured on rankings, impressions, clicks and CTR,
sessions, keyword coverage, indexed pages, referring domains, authority scores,
bounce\ix{bounce} and dwell, conversions and Core Web Vitals. The audit runs nine steps from
crawl and indexation through technical, on-page, content and backlink review to
competitor benchmarking, impact-versus-effort prioritisation, implementation and
re-measurement. Because search engines have evolved to mimic human preference,
good usability is now a requirement of good SEO rather than a separate concern.
\end{summarybox}

%%%%%%%%%%%%%%%%%%%%%%%%%%%%%%%%%%%%%%%%%%%%%%%%%%%%%%%%%%%%%%%%%%%%%%%%%%%%%%
%  Unit II appendices
%%%%%%%%%%%%%%%%%%%%%%%%%%%%%%%%%%%%%%%%%%%%%%%%%%%%%%%%%%%%%%%%%%%%%%%%%%%%%%

\startappendices

\chapter{Syllabus-to-Chapter Concordance}

This appendix exists so that coverage can be \emph{audited} rather than assumed.
The left-hand column reproduces the Unit~II syllabus descriptor phrase by phrase,
in the order the official document prints it. The right-hand column gives the
chapter and section that carries it. Every phrase resolves to a location, and no
numbered chapter of this volume contains material that is not traceable to a
phrase in this table.

\section*{First half --- Social Media Marketing}
\addcontentsline{toc}{section}{First half --- Social Media Marketing}

\begin{mmlongtable}{Concordance: Social Media Marketing}{tab:concordanceA2}%
{@{}P{0.40\textwidth}P{0.53\textwidth}@{}}
\rowcolor{ocre}\thd{Syllabus phrase} & \thd{Where it is covered}\\
\midrule
\endfirsthead
\rowcolor{ocre}\thd{Syllabus phrase} & \thd{Where it is covered}\\
\midrule
\endhead
\rlab{Social Media Marketing: Introduction} & Chapter~1 entire. Definition and
primary goal \S\,1.1; the trust asymmetry \S\,1.2; the omnichannel gearbox
\S\,1.3; myths \S\,1.4; marketer characteristics \S\,1.5; the customer lifecycle
orbit \S\,1.6.\\
\rowcolor{tablealt}
\rlab{Major Social Media Platforms for Marketing} & Chapter~2 entire. The
deployment matrix \S\,2.1; the Facebook algorithm filter \S\,2.2; X tactics and
hashtags \S\,2.3; the LinkedIn architecture \S\,2.4; visual content \S\,2.5;
platform choice and the psychology matrix \S\,2.6.\\
\rlab{Developing Data-driven Audience \& Campaign Insights} & Chapter~3 entire.
Definition \S\,3.1; the eight-step process \S\,3.2; the persona canvas \S\,3.3;
the seven analytic layers \S\,3.4; the content lifecycle matrix \S\,3.5; the
4-1-1 rule \S\,3.6; the 4~Rs \S\,3.7.\\
\rowcolor{tablealt}
\rlab{Social Media for Business} & Chapter~4 entire. Benefits \S\,4.1; the
ten-step plan \S\,4.2; tools \S\,4.3; blogs, podcasts and webinars \S\,4.4;
monitoring \S\,4.5.\\
\rlab{Creation \& Optimization of Social Media Campaigns} & Chapter~5 entire.
Required elements \S\,5.1; the campaign build \S\,5.2; optimisation levers
\S\,5.3; metric families \S\,5.4; funnel deployment \S\,5.5.\\
\end{mmlongtable}

\section*{Second half --- Search Engine Optimization}
\addcontentsline{toc}{section}{Second half --- Search Engine Optimization}

\begin{mmlongtable}{Concordance: Search Engine Optimization}{tab:concordanceB2}%
{@{}P{0.40\textwidth}P{0.53\textwidth}@{}}
\rowcolor{ocre}\thd{Syllabus phrase} & \thd{Where it is covered}\\
\midrule
\endfirsthead
\rowcolor{ocre}\thd{Syllabus phrase} & \thd{Where it is covered}\\
\midrule
\endhead
\rlab{Search Engine Optimization Fundamentals} & Chapter~6 entire. Definition and
objective \S\,6.1; why SEO matters \S\,6.2; SEO, PPC and SEM \S\,6.3; the SERP
\S\,6.4; intent and E-E-A-T \S\,6.5; the six-step process \S\,6.6; goals and
white hat against black hat \S\,6.7.\\
\rowcolor{tablealt}
\rlab{Keywords and SEO Content Plan} & Chapter~7 entire. Why research matters
\S\,7.1; keyword types \S\,7.2; the sweet spot \S\,7.3; the ten-step content
plan and cannibalisation \S\,7.4.\\
\rlab{SEO \& Business Objectives} & Chapter~8 entire. Objective-to-KPI alignment
\S\,8.1; integrating SEO into development \S\,8.2.\\
\rowcolor{tablealt}
\rlab{Writing SEO Content} & Chapter~9 entire. Page anatomy \S\,9.1; integrating
keywords naturally \S\,9.2; key considerations \S\,9.3.\\
\rlab{On-site \& off-site SEO} & Chapter~10 entire. On-page elements \S\,10.1;
technical SEO \S\,10.2; off-page SEO \S\,10.3; the comparison \S\,10.4; the
House of SEO \S\,10.5.\\
\rowcolor{tablealt}
\rlab{Optimize Organic Search Ranking} & Chapter~11 entire. Purpose \S\,11.1;
metrics \S\,11.2; the SEO audit \S\,11.3; the human--algorithm symbiosis
\S\,11.4; tools \S\,11.5.\\
\end{mmlongtable}

\begin{keybox}[Reading the ``etc.'' in the Syllabus]
The official descriptor ends several clauses with ``etc.''. That word is treated
in this series as licence to complete a named framework, not as licence to add
new topics. Where the syllabus says ``Creation \& Optimization of Social Media
Campaigns, etc.'', the chapter covers campaign creation and optimisation
completely rather than introducing an unnamed sixth topic.
\end{keybox}

\chapter{Previous-Year Question Bank --- Unit II}

Every question in this appendix is reproduced from an actual RVCE Marketing
Management Continuous Internal Evaluation or Semester End Examination paper for
this course. Answers and pointers follow the official schemes of valuation where
those were released.

\section*{Part A --- Two-Mark Questions}
\addcontentsline{toc}{section}{Part A --- Two-Mark Questions}

\begin{enumerate}[leftmargin=2.2em,itemsep=1.1ex]

\item \textbf{What is meant by data-driven audience insight in Social Media
      Marketing?} \pyqr{CIE-II, 26 May 2026 --- 2 M}\\
      The process of collecting and analysing user data --- demographics,
      interests, online behaviour, engagement patterns and preferences --- from
      social media platforms to understand the target audience better. It helps
      marketers create personalised campaigns and improve customer engagement.
      \hfill\emph{See} \S\,3.1.

\item \textbf{List any four elements required for creating an effective social
      media campaign.} \pyqr{CIE-II, 26 May 2026 --- 2 M}\\
      Clear campaign objectives; a defined target audience; engaging and creative
      content; performance measurement and analytics. \hfill\emph{See} \S\,5.1.

\item \textbf{Define campaign optimization in Social Media Marketing and mention
      its purpose.} \pyqr{CIE-II, 26 May 2026 --- 2 M}\\
      The process of improving social media campaign performance by analysing
      metrics and adjusting content, targeting and advertisements. Purpose:
      increase engagement, improve conversion rates, reduce marketing cost and
      maximise ROI. \hfill\emph{See} \S\,5.3.

\item \textbf{List any four benefits of using social media platforms for business
      marketing and customer engagement.} \pyqr{CIE-II, 26 May 2026 --- 2 M}\\
      Increased brand awareness; better customer interaction; cost-effective
      promotion; improved customer feedback collection. \hfill\emph{See} \S\,4.1.

\item \textbf{What is the primary goal of social media marketing?}
      \pyqr{SEE, Jun/Jul 2025 --- 2 M}\\
      To build brand awareness and engagement with a target audience on social
      platforms, and to convert that engagement into traffic, leads, sales and
      long-term customer loyalty. \hfill\emph{See} \S\,1.1.

\item \textbf{Define Hashtag. List the importance of Hashtag in social media
      marketing.} \pyqr{SEE, Oct/Nov 2023 --- 2 M}\\
      A hashtag is a word or phrase prefixed with \texttt{\#} that becomes a
      clickable, searchable label grouping all posts on that topic. Importance:
      increases discoverability beyond existing followers; groups a campaign
      under one trackable label; enables participation in trending
      conversations; collects user-generated content; and provides a measurable
      unit for campaign reporting. \hfill\emph{See} \S\,2.3.

\item \textbf{Why is keyword research important in SEO? / What is keyword
      research in SEO?}
      \pyqr{CIE-II, 7 May 2025; SEE, Oct/Nov 2023 --- 2 M}\\
      Keyword research is the process of identifying the phrases that real users
      type into search engines, together with their search volume, competition
      and intent. It matters because it ensures content targets phrases that
      people actually search, that the site can realistically rank for, and that
      carry commercial intent --- so effort produces revenue rather than vanity
      traffic. \hfill\emph{See} \S\,7.1.

\item \textbf{Differentiate between On-site SEO and Off-site SEO.}
      \pyqr{CIE-II, 7 May 2025 --- 2 M}\\
      On-site SEO is work done within your own website --- title tags, meta
      descriptions, headings, content quality, keyword placement, internal
      linking, site speed --- controlling how content is presented to search
      engines. Off-site SEO is work done outside your website --- backlinks,
      digital PR, social sharing, reviews, listings --- building the site's
      authority and reputation as judged by other sites.
      \hfill\emph{See} \S\,10.4.

\item \textbf{What is the purpose of optimizing organic search rankings in SEO?}
      \pyqr{CIE-II, 7 May 2025 --- 2 M}\\
      To appear higher in unpaid results for commercially relevant queries,
      capturing sustainable high-intent traffic at near-zero marginal cost,
      reducing dependence on paid advertising and gaining the credibility of an
      organic listing. \hfill\emph{See} \S\,11.1.

\item \textbf{What is the primary objective of Search Engine Optimization?}
      \pyqr{SEE, Jun/Jul 2025 --- 2 M}\\
      To increase the quantity and quality of organic traffic by improving the
      site's visibility and ranking for relevant queries, and to ensure the site
      satisfies the intent of the visitors it attracts. \hfill\emph{See} \S\,6.1.

\item \textbf{Why is it important to integrate SEO in website development?}
      \pyqr{SEE, Jun/Jul 2025 --- 2 M}\\
      Because architecture, URLs, rendering, speed, heading structure and
      tracking are all determined at build time. Retro-fitting SEO later is
      expensive, forces URL changes that risk losing accumulated authority, and
      often leaves staging \texttt{noindex} directives in production.
      \hfill\emph{See} \S\,8.2.

\item \textbf{Name two technical aspects of a website that can significantly
      impact its organic search ranking.}
      \pyqr{CIE-II Quiz, 24 Jul 2024 --- 2 M}\\
      Site loading speed and Core Web Vitals, and mobile responsiveness. Also
      acceptable: HTTPS, crawlability and indexability, XML sitemap, URL
      structure. \hfill\emph{See} \S\,10.2.

\item \textbf{Define SEO and explain why it is important for online businesses.}
      \pyqr{CIE-II Quiz, 24 Jul 2024 --- 2 M}\\
      SEO is the set of techniques used to improve a website's visibility and
      ranking in the organic results of search engines. It matters to online
      businesses because organic search delivers about 53\% of all website
      traffic, carries the highest buyer intent of any channel, costs nothing per
      click, compounds over time, and confers more credibility than paid
      listings. \hfill\emph{See} \S\,6.1 and \S\,6.2.

\item \textbf{What are two key considerations when writing SEO-friendly content
      for web pages?} \pyqr{CIE-II Quiz, 24 Jul 2024 --- 2 M}\\
      (1) Match the search intent --- the format and depth must answer what the
      user actually wants. (2) Integrate keywords naturally in the title, H1,
      opening paragraph and subheadings using semantic variants, at roughly
      3--5\% density, without stuffing. Also acceptable: a unique title tag and
      meta description; a scannable structure. \hfill\emph{See} \S\,9.3.

\end{enumerate}

\section*{Part B --- Eight- and Ten-Mark Questions}
\addcontentsline{toc}{section}{Part B --- Eight- and Ten-Mark Questions}

\begin{enumerate}[leftmargin=2.2em,itemsep=1.4ex]

\item \textbf{Analyse how the Facebook Algorithm Filter evaluates content
      distribution. What specific practices are penalized (filtered out) and
      which ones are rewarded (propelled through)?}
      \pyqr{CIE-II, 26 May 2026 --- 10 M, both sets}\\
      Structure the answer as: (i) what the algorithm is and why it exists;
      (ii) the five evaluation criteria --- engagement history, relevance,
      content type preference, timeliness, quality and originality; (iii) the
      penalised list --- clickbait headlines, repeated sources posted
      back-to-back, engagement bait, fake news, spam, low-quality copied
      content, excessive promotion; (iv) the rewarded list --- time spent on the
      post after full load, video actions (unmuting, full-screen, HD), Facebook
      Live at 300\% more watch time than recorded video, original content,
      meaningful interactions, informative posts, consistent posting.
      \hfill\emph{See} \S\,2.2.

\item \textbf{Based on ``The Content Lifecycle Matrix,'' how should a marketer
      adjust the objective, tone, and format of their content as a prospect
      transitions from the Early-Stage to the Late-Stage?}
      \pyqr{CIE-II, 26 May 2026 --- 10 M}\\
      Reproduce the three-by-three matrix and then explain the trajectory: the
      objective moves from attention to trust to action; the tone moves from fun
      and entertaining to educational to informative and conversion-focused; the
      format moves from curated news, tips and visual content to contests,
      newsletter subscriptions and event invitations, and finally to
      click-to-purchase, demonstrations and trial sign-ups.
      \hfill\emph{See} \S\,3.5.

\item \textbf{Explain the stages of ``The Customer Lifecycle Orbit'' and discuss
      the statistical differences in selling to new prospects versus existing
      customers.} \pyqr{CIE-II, 26 May 2026 --- 10 M}\\
      Name and explain the six stages, then give the statistics: the probability
      of selling to a new prospect is roughly 5--20\%, against 60--70\% for an
      existing customer; repeat customers spend up to 67\% more; 49\% of
      marketers achieve higher ROI by focusing on engagement. Conclude with the
      implication --- retention is structurally cheaper and more profitable than
      acquisition, so social content must serve the whole orbit, not just
      awareness. \hfill\emph{See} \S\,1.6.

\item \textbf{Detail the ``4 Rs of Content Scaling.'' How does this strategic
      framework help marketers maximize the use of their content assets?}
      \pyqr{CIE-II, 26 May 2026 --- 10 M}\\
      Give Reorganize, Rewrite, Retire and Redesign with an example each, note
      the alternative Repurpose, Refresh, Repost and Recombine formulation, and
      close with the benefits --- saves time and cost, increases reach, improves
      SEO performance, enhances engagement and maximises ROI from content
      assets. \hfill\emph{See} \S\,3.7.

\item \textbf{Explain the importance of social media marketing for businesses
      today. How can businesses integrate social media platforms into their
      overall marketing strategy?} \pyqr{CIE-II, 7 May 2025 --- 10 M}\\
      \emph{Importance} --- the trust asymmetry (90\% believe peers, 48\% trust
      businesses), reach, cost-effectiveness, two-way engagement, feedback,
      social selling and SEO fuel. \emph{Integration} --- the omnichannel
      gearbox, the ten-step plan, mapping content to the customer lifecycle
      orbit, and synchronising social with website, e-mail, search and offline
      channels so the customer experiences one continuous conversation.
      \hfill\emph{See} \S\,1.2, \S\,1.3 and \S\,4.2.

\item \textbf{Compare the roles of Facebook, Instagram, LinkedIn, and Twitter in
      social media marketing. How should a business choose the right platform?}
      \pyqr{CIE-II, 7 May 2025 --- 10 M}\\
      Use the platform deployment matrix covering core mechanic, best use and key
      features for each, then the six selection criteria, closing with the
      platform psychology matrix --- each platform serves a different
      psychological function in the decision. \hfill\emph{See} \S\,2.1 and
      \S\,2.6.

\item \textbf{Explain the importance of data-driven insights in social media
      marketing. Describe the process of developing audience insights using
      social media analytics tools. How can these insights inform the
      development and optimization of social media campaigns?}
      \pyqr{CIE-I, 1 Jul 2024 --- 10 M}\\
      Cover the definition, the eight-step insight process, the tools, and then
      the application --- targeting, creative, tone, format, schedule and budget
      allocation, plus in-flight optimisation using the eight levers.
      \hfill\emph{See} Chapter~3 and \S\,5.3.

\item \textbf{Illustrate how social media analytics contribute to creating
      audience profiles.} \pyqr{SEE, Jun/Jul 2025, Q4a --- 8 M}\\
      The seven layers: demographic, temporal, interest, behavioural,
      conversational, comparative and conversion.
      \hfill\emph{See} \S\,3.4.

\item \textbf{Analyse how businesses can utilize Facebook and Instagram for
      different stages of a marketing funnel.}
      \pyqr{SEE, Jun/Jul 2025, Q3a --- 8 M}\\
      Cover awareness, consideration, conversion and retention or advocacy on
      both platforms. \hfill\emph{See} \S\,5.5.

\item \textbf{Discuss the key metrics used to measure the success of a social
      media marketing campaign and explain the importance of it for evaluating
      performance.} \pyqr{SEE, Oct/Nov 2023, Q3a --- 8 M}\\
      The six metric families and the six reasons measurement matters.
      \hfill\emph{See} \S\,5.4.

\item \textbf{Analyse the differences between on-site and off-site SEO
      techniques. Provide examples for each.}
      \pyqr{CIE-II, 7 May 2025 --- 10 M}\\
      Use the comparison table and the House of SEO, with the worked examples
      given there. \hfill\emph{See} \S\,10.4 and \S\,10.5.

\item \textbf{Explain the key components of Search Engine Optimization,
      including on-page, off-page, and technical SEO. How do these elements
      contribute to improving website ranking on search engines?}
      \pyqr{CIE-I, 16 Apr 2026 --- 10 M}\\
      \emph{On-page SEO} (content optimisation) --- keyword optimisation in
      titles, headings and content; meta tags; URL structure; internal linking;
      high-quality relevant content. \emph{Off-page SEO} (external factors) ---
      backlinks from other websites; social media sharing; influencer
      collaborations; online reviews. \emph{Technical SEO} (backend
      optimisation) --- website speed; mobile friendliness; secure HTTPS; XML
      sitemap; crawlability and indexing. Then explain the contribution:
      technical SEO makes the site discoverable, on-page SEO makes it relevant,
      and off-page SEO makes it authoritative --- and rankings are a function of
      all three together. \hfill\emph{See} Chapter~10.

\item \textbf{What are the key metrics used to measure organic search ranking?
      How can businesses monitor and improve these metrics over time? Describe
      the process of conducting an SEO audit and implementing changes to optimize
      organic search ranking.} \pyqr{CIE-II, 24 Jul 2024 --- 10 M}\\
      The metric table, then the nine-step audit process, closing with the
      six-step SEO process as the ongoing improvement cycle.
      \hfill\emph{See} \S\,11.2, \S\,11.3 and \S\,6.6.

\item \textbf{An educational institution is struggling to rank for competitive
      keywords related to online courses. Conduct an SEO audit of their website,
      identify key issues affecting organic search ranking, and recommend
      actionable solutions.} \pyqr{CIE-II, 24 Jul 2024 --- 10 M}\\
      The worked diagnostic is set out in full in the chapter.
      \hfill\emph{See} \S\,11.3.

\item \textbf{You are writing a blog post about a complex technical topic.
      Describe how you would integrate relevant keywords naturally into the
      content while ensuring it remains engaging and informative for readers.}
      \pyqr{CIE-II, 24 Jul 2024 --- 10 M}\\
      The nine-point method, anchored on ``write for the reader first, then
      optimise'': keyword placement in high-value positions, semantic variants,
      3--5\% density, People-Also-Ask coverage, jargon explanation, visual
      breaks, internal linking and a read-aloud check.
      \hfill\emph{See} \S\,9.2.

\item \textbf{Develop an SEO content plan for an e-commerce startup targeting
      eco-friendly household products.} \pyqr{SEE, Jun/Jul 2025, Q3b --- 8 M}\\
      \emph{Objective:} organic revenue from high-intent shoppers within nine
      months. \emph{Keyword research:} cluster into topics --- eco-friendly
      cleaning, plastic-free kitchen, sustainable laundry, bamboo and
      natural-fibre goods --- targeting long-tail commercial phrases such as
      ``biodegradable dishwashing liquid India'' (volume 250--1,000+, difficulty
      under 40, clear buyer intent). \emph{Architecture:} one pillar page per
      cluster, supported by product and comparison pages, all internally linked.
      \emph{Content by intent:} informational blogs (``how harmful are
      conventional detergents?''), commercial comparisons (``best plastic-free
      kitchen swaps''), and transactional category and product pages with unique
      copy. \emph{On-page:} unique title tags under 70 characters, 155--160
      character meta descriptions, one H1 per page, ALT text on all product
      images, 3--5\% keyword density, breadcrumb navigation.
      \emph{Technical:} a fast mobile-first storefront, HTTPS, XML sitemap,
      Product and Review schema for rich results, canonical tags on filtered
      category URLs. \emph{Off-page:} digital PR on plastic-waste data, guest
      posts on sustainability blogs, Google Business Profile, review generation.
      \emph{Measurement:} keyword coverage, organic sessions, organic conversion
      rate and revenue, reviewed monthly. \hfill\emph{See} \S\,7.4.

\item \textbf{Construct a targeted marketing plan using SEO tools like Google
      Analytics and SEMrush for a B2B company.}
      \pyqr{SEE, Jun/Jul 2025, Q4b --- 8 M}\\
      \emph{Step 1 --- baseline in Google Analytics 4:} identify current organic
      sessions, top landing pages, acquisition channels, and which pages assist
      conversions. \emph{Step 2 --- opportunity analysis in SEMrush:} run
      Organic Research on the top five competitors, use Keyword Gap to find
      phrases they rank for and you do not, and the Keyword Magic Tool to build
      clusters around the buyer's problems. \emph{Step 3 --- persona and intent
      mapping:} align keywords to the four to six B2B personas (Executive
      Sponsor, Decision Maker, End User), matching informational content to
      early-stage roles and commercial content to decision makers.
      \emph{Step 4 --- content plan:} pillar pages per solution area, supported
      by case studies, ROI calculators, whitepapers and comparison pages.
      \emph{Step 5 --- technical and on-page fixes} from the SEMrush Site Audit,
      prioritised by impact. \emph{Step 6 --- off-page:} industry publication
      guest posts, original research for digital PR, LinkedIn thought
      leadership. \emph{Step 7 --- measurement:} SEMrush Position Tracking for
      rankings and share of voice; GA4 goals and UTM-tagged campaigns for
      marketing-qualified leads, cost per lead and pipeline influence; monthly
      review and reallocation. \hfill\emph{See} \S\,11.5 and \S\,8.1.

\item \textbf{Explain the concept of Search Engine optimization and its
      significance in digital marketing.}
      \pyqr{SEE, Oct/Nov 2023, Q4a --- 8 M}\\
      Define SEO, explain the SEO, PPC and SEM distinction, the six-step process,
      and the three pillars --- technical, on-page and off-page.
      \emph{Significance:} 53\% of all traffic, the highest buyer intent,
      compounding at near-zero marginal cost, greater credibility than
      advertisements, competitive defence, and measurable results.
      \hfill\emph{See} Chapter~6.

\item \textbf{Describe the role of content in SEO. How can business create high
      quality SEO friendly content?}
      \pyqr{SEE, Oct/Nov 2023, Q4b --- 8 M}\\
      \emph{Role:} content is king --- search engines store and rank text, so
      high-volume, high-quality relevant content both gives users a reason to
      stay and return, and gives engines the material from which rankings are
      computed. It is also the asset that earns backlinks and satisfies E-E-A-T.
      \emph{How:} research keywords and intent first; write for the reader and
      optimise afterwards; use the anatomy of a top-ranking page; maintain
      3--5\% keyword density with semantic variants; write unique title tags and
      meta descriptions; structure with one H1 and logical H2 and H3 headings;
      add rich media for dwell time; link internally with descriptive anchors;
      keep content fresh through historical audits; and ensure it is more
      comprehensive and more readable than what currently ranks.
      \hfill\emph{See} Chapter~9.

\item \textbf{Enumerate briefly on the best practices for creating compelling and
      shareable content on social media platforms to increase brand visibility
      and engagement.} \pyqr{SEE, Oct/Nov 2023, Q3b --- 8 M}\\
      Apply the seven-point quality gate from Volume~1; lead with visuals ---
      40 times more shareable, 94\% more views with relevant images, and 135\%
      greater organic reach for video; follow the 4-1-1 rule; write natively for
      each platform rather than cross-posting; hook within the first three
      seconds; use one or two hashtags only; post consistently at
      audience-active times; invite participation through polls, questions and
      contests; scale one master asset with the 4~Rs; and amplify through
      communities, employees and influencers.
      \hfill\emph{See} \S\,2.5, \S\,3.6 and \S\,3.7.

\end{enumerate}

\section*{How Unit II Is Examined}
\addcontentsline{toc}{section}{How Unit II Is Examined}

\begin{keybox}[The Pattern Across Papers]
In the Semester End Examination, Unit~II is Questions~3 and~4 with internal
choice --- 16 marks in two eight-mark sub-divisions --- plus two or three Part~A
objective questions. The second Continuous Internal Evaluation test is typically
dominated by this unit.

\smallskip
The highest-probability ten-mark topics are the \textbf{Facebook algorithm
filter} (asked twice in one paper), the \textbf{content lifecycle matrix}, the
\textbf{4~Rs of content scaling}, the \textbf{customer lifecycle orbit} with its
statistics, \textbf{on-site against off-site SEO}, and the \textbf{SEO audit
process}.

\smallskip
Note that the Semester End Examination also mixes social media marketing and SEO
into applied ``design a plan for X'' questions. Always answer those with a named
framework, not free-form opinion.
\end{keybox}

\chapter{Supplementary Topic: Search Everywhere Optimization}

\begin{warnbox}[What This Appendix Is, and Why It Is Not a Chapter]
Search Everywhere Optimization --- sometimes called SEO 2.0 --- appears in the
supplementary lecture material for this unit but is named by no unit descriptor
in the official syllabus. It is placed in an appendix so that the rule governing
this series holds: the numbered chapters carry the syllabus and nothing else. It
is included rather than dropped because it is directly useful for any question
that asks about \emph{new} digital channels or about why traffic can rise while
conversions stay flat.
\end{warnbox}

\section{The Scale of Search Outside Google}

\begin{keybox}[The Figures]
Around \textbf{13.7 billion} searches happen daily across the internet, and
\textbf{73\%} of all searches happen \emph{outside} Google. Most marketers are
still playing a Google-only game --- which means a business winning at Google can
simultaneously be losing customers.
\end{keybox}

\section{The Symptom and the Disease}

\begin{description}
\item[The symptom] Traffic looks decent but conversions are flat. The business is
showing up for the search and missing the decision.
\item[The disease] Optimising for visibility in one place while the customer is
deciding everywhere else.
\end{description}

\section{The Death of the Funnel}

The modern journey is not funnel-shaped and linear. It is a \emph{decision
constellation} of micro-choices happening in different places at the same time.

\begin{figure}[htbp]
\centering
\begin{tikzpicture}[font=\sffamily\footnotesize,
    nd/.style={draw=ocre!70,line width=0.8pt,fill=ocrepale,rounded corners=3pt,
               align=center,text width=2.05cm,inner sep=4pt,
               font=\sffamily\scriptsize,text=slate}]
  \node[draw=slate,line width=1.1pt,fill=white,circle,minimum size=2.3cm,
        align=center,font=\sffamily\bfseries\scriptsize,text=slate] (c) at (0,0)
        {THE\\ DECISION};
  \foreach \ang/\n/\lbl in {%
      90/a/{\textbf{Click}\\ search engine},
      45/b/{\textbf{Trust}\\ forums},
      0/c2/{\textbf{Buy}\\ marketplaces},
     -45/d/{\textbf{Try}\\ app stores},
     -90/e/{\textbf{Think}\\ video, podcasts},
    -135/f/{\textbf{Believe}\\ AI assistants},
     180/g/{\textbf{Follow}\\ social feeds},
     135/h/{\textbf{Better}\\ video reviews}}{
     \node[nd] (\n) at (\ang:3.75cm) {\lbl};
     \draw[ocre!70,line width=0.7pt,dashed] (\n) -- (c);}
\end{tikzpicture}
\caption[The decision constellation]%
        {The decision constellation. Because these micro-decisions happen in
         parallel rather than in sequence, optimising a single channel optimises
         a single point in a distributed process.}
\label{fig:constellation}
\end{figure}

\section{Visibility Against Validation}

\begin{keybox}[The Distinction]
\textbf{Visibility} is what you do --- showing up in the search engine, creating
platform-specific posts. \textbf{Validation} is what others say about what you do
--- organic mentions in conversations, authentic social proof.

\smallskip
\textbf{Visibility without validation is just noise.} You need both to secure the
decision.
\end{keybox}

\section{Prioritisation: the RICE Framework}

You do not have to be everywhere. You need a prioritisation strategy.

\begin{mmlongtable}{The RICE prioritisation framework}{tab:rice}%
{@{}P{0.16\textwidth}P{0.77\textwidth}@{}}
\rowcolor{ocre}\thd{Criterion} & \thd{The question it asks}\\
\midrule
\endfirsthead
\rowcolor{ocre}\thd{Criterion} & \thd{The question it asks}\\
\midrule
\endhead
\rlab{Reach} & How many users in your specific audience will this
micro-decision impact?\\
\rowcolor{tablealt}
\rlab{Impact} & How deeply does this platform influence the final purchase
decision?\\
\rlab{Confidence} & How certain are you that your brand can execute
authentically here?\\
\rowcolor{tablealt}
\rlab{Ease} & What operational effort is required to maintain a strategic
presence?\\
\end{mmlongtable}

\chapter{Glossary of Key Terms}

\begin{mmlongtable}{Key terms for rapid revision}{tab:glossary2}%
{@{}P{0.27\textwidth}P{0.66\textwidth}@{}}
\rowcolor{ocre}\thd{Term} & \thd{One-line meaning}\\
\midrule
\endfirsthead
\rowcolor{ocre}\thd{Term} & \thd{One-line meaning}\\
\midrule
\endhead
\rlab{4-1-1 rule} & Four value posts, one soft promotion, one hard promotion
(Joe Pulizzi)\\
\rowcolor{tablealt}
\rlab{4 Rs of content scaling} & Reorganize, Rewrite, Retire, Redesign\\
\rlab{Backlink} & An inbound link from another site --- the currency of off-page
SEO\\
\rowcolor{tablealt}
\rlab{Campaign optimization} & Improving campaign performance by analysing
metrics and adjusting content, targeting and advertisements\\
\rlab{Content lifecycle matrix} & Early, mid and late-stage mapping of objective,
tone and format\\
\rowcolor{tablealt}
\rlab{Core Web Vitals} & LCP, INP and CLS --- the technical experience signals
that feed into ranking\\
\rlab{Customer lifecycle orbit} & Awareness, engagement, purchase, retention and
loyalty, growth, advocacy\\
\rowcolor{tablealt}
\rlab{Data-driven audience insight} & Collecting and analysing platform user data
to understand the audience\\
\rlab{E-E-A-T} & Experience, Expertise, Authoritativeness, Trustworthiness\\
\rowcolor{tablealt}
\rlab{Engagement bait} & Explicitly asking for likes, shares or comments ---
penalised by the algorithm\\
\rlab{Facebook algorithm filter} & The ranking system deciding which posts reach
the news feed\\
\rowcolor{tablealt}
\rlab{Featured snippet} & ``Position zero'' --- a direct answer extracted onto
the SERP\\
\rlab{Hashtag} & A \texttt{\#}-prefixed searchable label grouping posts on a
topic\\
\rowcolor{tablealt}
\rlab{Head keyword} & A one- or two-word phrase: huge volume, brutal
competition, vague intent\\
\rlab{Keyword cannibalisation} & Two of your own pages competing for the same
keyword\\
\rowcolor{tablealt}
\rlab{Keyword density} & Share of the content occupied by the focus keyword;
industry norm 3--5\%\\
\rlab{Link gravity} & The principle that link quality and relevance outweigh raw
link count\\
\rowcolor{tablealt}
\rlab{Long-tail keyword} & A four-or-more-word phrase: low volume, low
competition, high intent\\
\rlab{Off-page SEO} & A site's authority as determined by what other sites say
about it\\
\rowcolor{tablealt}
\rlab{On-page SEO} & How well a site's own content and code are presented to
search engines\\
\rlab{Panda / Penguin} & Google updates penalising thin content and manipulative
link-building\\
\rowcolor{tablealt}
\rlab{RankBrain / Hummingbird} & Google algorithms, from 2015, focused on user
intent rather than keyword strings\\
\rlab{RICE} & Reach, Impact, Confidence, Ease --- a channel prioritisation
framework\\
\rowcolor{tablealt}
\rlab{SEM} & Umbrella term covering SEO plus PPC\\
\rlab{SEO} & Improving organic visibility and ranking in search engines\\
\rowcolor{tablealt}
\rlab{SERP} & Search engine results page\\
\rlab{Share of voice} & Conversations about your brand against conversations
about competitors\\
\rowcolor{tablealt}
\rlab{Social media marketing} & Creating and sharing content on social networks
to achieve marketing and branding goals\\
\rlab{Social media monitoring} & Tracking mentions of brand, competitors,
keywords and hashtags\\
\rowcolor{tablealt}
\rlab{Technical SEO} & The crawlable, indexable, fast, secure infrastructure of
the site\\
\rlab{Trust asymmetry} & Roughly 90\% believe peers against about 48\% who trust
businesses directly\\
\rowcolor{tablealt}
\rlab{White hat / black hat SEO} & Techniques within search-engine guidelines,
against techniques that violate them\\
\end{mmlongtable}

\chapter{Framework Quick-Reference Sheet}

Everything in this volume that is a numbered list, in one place, for the night
before the examination.

\begin{mmlongtable}{Framework quick reference}{tab:quickref2}%
{@{}P{0.28\textwidth}P{0.65\textwidth}@{}}
\rowcolor{ocre}\thd{Framework} & \thd{The sequence}\\
\midrule
\endfirsthead
\rowcolor{ocre}\thd{Framework} & \thd{The sequence}\\
\midrule
\endhead
\rlab{Trust asymmetry} & About 90\% believe peers; about 48\% trust businesses;
word of mouth valued about 41\% above social recommendations\\
\rowcolor{tablealt}
\rlab{Omnichannel gearbox} & Social selling, SEO fuel, omnichannel
synchronisation\\
\rlab{Customer lifecycle orbit} & Awareness \ra{} engagement \ra{} purchase
\ra{} retention and loyalty \ra{} growth \ra{} advocacy\\
\rowcolor{tablealt}
\rlab{Retention statistics} & New prospect under 20\% (5--20\%); existing
customer over 60\% (60--70\%); repeat customers spend up to 67\% more; 49\%
report higher ROI from engagement\\
\rlab{Eight-step insight process} & Define the question \ra{} collect \ra{}
consolidate \ra{} segment \ra{} find patterns \ra{} build personas \ra{} apply
\ra{} iterate\\
\rowcolor{tablealt}
\rlab{Seven analytic layers} & Demographic, temporal, interest, behavioural,
conversational, comparative, conversion\\
\rlab{Persona canvas} & Background and goals; information sources; messaging and
topics; purchase role and objections\\
\rowcolor{tablealt}
\rlab{Content lifecycle matrix} & Objective: attention \ra{} trust \ra{} action.
Tone: fun \ra{} educational \ra{} decisive. Format: curated \ra{} participatory
\ra{} transactional\\
\rlab{4-1-1 rule} & Four value + one soft promotion + one hard promotion per six
posts\\
\rowcolor{tablealt}
\rlab{4 Rs of content scaling} & Reorganize, Rewrite, Retire, Redesign --- also
described as Repurpose, Refresh, Repost, Recombine\\
\rlab{Facebook filter --- rewarded} & Time spent, video actions, Live, original
content, meaningful interaction, informative posts, consistency\\
\rowcolor{tablealt}
\rlab{Facebook filter --- penalised} & Clickbait, repeated sources, engagement
bait, fake news, spam, copied content, excessive promotion\\
\rlab{Platform choice criteria} & Persona location, B2C against B2B, format
capability, funnel stage, resources, measured ROI\\
\rowcolor{tablealt}
\rlab{Platform psychology} & Discovery, education, verification, purchasing,
aspiration, validation, the click\\
\rlab{Ten-step SMM plan} & Goals \ra{} audit \ra{} audience \ra{} platforms
\ra{} content strategy \ra{} calendar \ra{} resources \ra{} publish and engage
\ra{} amplify \ra{} measure\\
\rowcolor{tablealt}
\rlab{Campaign elements} & Clear objectives, defined audience, engaging creative,
measurement\\
\rlab{Eight optimisation levers} & Creative, copy and CTA, audience, placement,
timing and frequency, budget, landing page, format\\
\rowcolor{tablealt}
\rlab{Six metric families} & Reach and volume, engagement, traffic, conversion,
sentiment and advocacy, efficiency\\
\rlab{SEO, PPC, SEM} & Organic, paid, and the umbrella over both\\
\rowcolor{tablealt}
\rlab{SERP result types} & Paid, organic, universal or vertical\\
\rlab{Six-step SEO process} & Keyword research \ra{} reporting and goals \ra{}
content building \ra{} page optimisation \ra{} social and link building \ra{}
follow-up analysis\\
\rowcolor{tablealt}
\rlab{Keyword intent types} & Informational, navigational, commercial
investigation, transactional\\
\rlab{Keyword lengths} & Head (1--2), body (2--3), long tail (4+)\\
\rowcolor{tablealt}
\rlab{Keyword sweet spot} & Volume 250--1,000+ $\times$ difficulty under 40
$\times$ commercial intent\\
\rlab{Ten-step content plan} & Objectives \ra{} keyword clusters \ra{} intent
mapping \ra{} content audit \ra{} gap analysis \ra{} pillar architecture \ra{}
one keyword per page \ra{} calendar \ra{} targets \ra{} monthly review\\
\rowcolor{tablealt}
\rlab{On-page five elements} & Page copy, title tags, meta data, heading tags,
interlinking\\
\rlab{On-page limits} & Title about 70 characters; meta description 155--160
characters; keyword density 3--5\%; one H1 per page\\
\rowcolor{tablealt}
\rlab{House of SEO} & Technical (foundation, discoverable), on-page (floor,
relevant), off-page (roof, authoritative)\\
\rlab{Nine-step SEO audit} & Crawl \ra{} indexation \ra{} technical \ra{}
on-page \ra{} content \ra{} backlinks \ra{} benchmark \ra{} prioritise \ra{}
implement and re-measure\\
\rowcolor{tablealt}
\rlab{RICE (Appendix~C)} & Reach, Impact, Confidence, Ease\\
\end{mmlongtable}


\backmatter
\printindex

\end{document}
